\exercisesection{}

Unless stated otherwise, the Banach spaces and Hilbert spaces below are assumed to be complex.

\begin{enumerate}
\def\labelenumi{\arabic{enumi})}
\tightlist
\item
  \begin{enumerate}
  \def\labelenumii{\alph{enumii})}
  \tightlist
  \item
    There exist vector spaces E, F, G and partial operators u, v from E to F, and w from F to G, such that \(w \circ (u + v) \neq w \circ u + w \circ v\) .
  \end{enumerate}
\end{enumerate}

\begin{enumerate}
\def\labelenumi{\alph{enumi})}
\setcounter{enumi}{1}
\item
  Give an example of a closable operator u on a topological vector space such that the domain of the closure of u is strictly contained in the closure of the domain of u.
\item
  Give an example of closed operators u and v on a Hilbert space such that u is continuous but \(u \circ v\) is not closed.
\end{enumerate}

\begin{enumerate}
\def\labelenumi{\arabic{enumi})}
\setcounter{enumi}{1}
\tightlist
\item
  Let E and F be Banach spaces, and let G be a subspace of E. Let u be a finite-rank linear map from G into F. Suppose that u is not continuous on G.
\end{enumerate}

\begin{enumerate}
\def\labelenumi{\alph{enumi})}
\item
  The partial operator from E to F with domain G defined by u is not closable.
\item
  Suppose that E and G are Hilbert spaces. Compute the adjoint of the partial operator defined by u.
\end{enumerate}

\begin{enumerate}
\def\labelenumi{\arabic{enumi})}
\setcounter{enumi}{2}
\tightlist
\item
  Let K = R or C. Let E, F be topological vector spaces over K and u an operator with dense domain from E into F. Let D' be the set of continuous linear forms \(\lambda \in F'\) such that there exists a continuous linear form \(\mu \in E'\) coinciding with \(\lambda \circ u\) on the domain of u.
\end{enumerate}

\begin{enumerate}
\def\labelenumi{\alph{enumi})}
\tightlist
\item
  The set D' is a vector subspace of F', the linear form \(\mu\) is determined in a unique way by \(\lambda\) , and the map \(\lambda \mapsto \mu\) from D' to E' is linear.
\end{enumerate}

One calls the transpose of u, and denotes \(^{t}u\) the partial operator from F' into E' whose domain is D' and which is given by \(\lambda\mapsto\mu\) for \(\lambda\in D'\) .

\begin{enumerate}
\def\labelenumi{\alph{enumi})}
\setcounter{enumi}{1}
\item
  Suppose that E and F are locally convex. The operator u is closable if and only if the duality \(\mathrm{F} \times \mathrm{dom}(^{t}u) \to \mathrm{K}\) is separating.
\item
  Suppose that E and F are locally convex and that F is semi-reflexive. The operator u is closable if and only if \(^{t}u\) has dense domain.
\end{enumerate}

\begin{enumerate}
\def\labelenumi{\arabic{enumi})}
\setcounter{enumi}{3}
\item
  Let E be a complex Hilbert space and u an injective self-adjoint partial operator on E. The partial operator \(u^{-1}\) is self-adjoint.
\item
  Let E be a complex Banach space. A continuous semigroup on E is called a map \(s\colon\mathbf{R}_{+}\to\mathcal{L}(\mathrm{E})\) such that
\end{enumerate}

\begin{enumerate}
\def\labelenumi{(\roman{enumi})}
\item
  \(s(0) = 1_{\mathrm{E}}\) ;
\item
  The map from \(R_{+} \times E \to E\) defined by \((t, x) \mapsto s(t)x\) is continuous;
\item
  For all \(t_1\) and \(t_2\) in \(\mathbf{R}_{+}\) , we have \(s(t_1 + t_2) = s(t_1)s(t_2)\) .
\end{enumerate}

We consider a continuous semigroup s on E.

\begin{enumerate}
\def\labelenumi{\alph{enumi})}
\tightlist
\item
  There exist real numbers C ≥ 0 and M ≥ 0 such that
\end{enumerate}

\[
\| s (t) \| \leqslant \mathrm{Ce} ^ {\mathrm{Mt}}
\]

for all \(t \in R_{+}\) . (Let a \textgreater{} 0; show that \(\|s(t)\|\) is bounded for \(t \in [0, a]\) .) We fix such numbers in the rest of the exercise.

\begin{enumerate}
\def\labelenumi{\alph{enumi})}
\setcounter{enumi}{1}
\item
  Let D be the set of \(x \in E\) such that the mapping \(\psi_{x}\) of \(R_{+}\) to E defined by \(t \mapsto s(t)x\) is differentiable at 0. The space D is dense in E; the map u from D into E defined by \(u(x) = -\psi_{x}'(0)\) is a closed partial operator on E, which is called the infinitesimal generator of the semigroup s.
\item
  Let \(\lambda\) a complex number such that \(\mathcal{R}(\lambda) < -\mathrm{M}\) . The formula
\end{enumerate}

\[
v _ {\lambda} = - \int_ {\mathbf {R} _ {+}} e ^ {\lambda t} s (t) d t
\]

defines an endomorphism of E.

\(d)\) The image of \(v_{\lambda}\) is contained in \(\operatorname{dom}(u)\) and we have \(u \circ v_{\lambda} = -1_{\mathrm{E}} - \lambda v_{\lambda}\) .

\begin{enumerate}
\def\labelenumi{\alph{enumi})}
\setcounter{enumi}{4}
\tightlist
\item
  One has \(u \circ v_{\lambda} = v_{\lambda} \circ u\) .
\end{enumerate}

\(f)\) The complex number \(\lambda\) belongs to the resolvent set of \(u\) and we have \(v_{\lambda} = -\mathrm{R}(u,\lambda)\) .

\begin{enumerate}
\def\labelenumi{\alph{enumi})}
\setcounter{enumi}{6}
\tightlist
\item
  One has
\end{enumerate}

\[
\| v _ {\lambda} ^ {n} \| \leqslant \frac {\mathrm{C}}{| \mathcal {R} (\lambda) + \mathrm{M} | ^ {n}}
\]

for all \(n \in N\) .

\begin{enumerate}
\def\labelenumi{\arabic{enumi})}
\setcounter{enumi}{5}
\tightlist
\item
  Let E be a complex Banach space and u a closed partial operator with dense domain on E. Suppose there exist real numbers \(C \geqslant 0\) and \(M \geqslant 0\) such that
\end{enumerate}

\begin{enumerate}
\def\labelenumi{(\roman{enumi})}
\item
  The set {]}−∞, −M{[} is contained in the resolvent set of u;
\item
  For every integer \(n \in \mathbf{N}\) and every real number \(\lambda < -\mathrm{M}\) , we have
\end{enumerate}

\[
\left\| \mathrm{R} (u, \lambda) ^ {n} \right\| \leqslant \mathrm{C} | \lambda + \mathrm{M} | ^ {- n}.
\]

\begin{enumerate}
\def\labelenumi{\alph{enumi})}
\tightlist
\item
  For every real number \(y < -\mathrm{M}\) , we denote
\end{enumerate}

\[
u _ {y} = - y 1 _ {\mathrm{E}} + y ^ {2} \mathrm{R} (u, y) \in \mathcal {L} (\mathrm{E}).
\]

The endomorphisms \(u_{y}\) commute pairwise; for every \(x \in \text{dom}(u)\) , we have \(u_{y}(x) \to u(x)\) as \(y \to -\infty\) .

\begin{enumerate}
\def\labelenumi{\alph{enumi})}
\setcounter{enumi}{1}
\item
  The map from \(R_{+}\) in \(\mathcal{L}(E)\) defined by \(s_{y}(t)=e^{-tu_{y}}\) is a continuous semigroup on E such that \(\|s_{y}(t)\|\leqslant Ce^{Mt}\) for all \(t\in R_{+}\) .
\item
  For all \(y_{1}, y_{2}\) and \(t\) in \(\mathbf{R}_{+}^{*}\) , and every \(x\) in \(\mathrm{dom}(u)\) , we have
\end{enumerate}

\[
\left\| s _ {y _ {1}} (t) x - s _ {y _ {2}} (t) x \right\| \leqslant t \mathrm{C} ^ {2} \left\| u _ {y _ {1}} (x) - u _ {y _ {2}} (x) \right\|.
\]

\(d)\) Let \(t \in \mathbf{R}_{+}\) . For every \(x\) in E, the function defined on \(\mathbf{R}_{+}^{*}\) by \(y \mapsto s_{y}(t)x\) converges as \(y \to -\infty\) to an element \(s(t)x\) of E. The map \(s(t)\) is linear and continuous with norm \(\leqslant \mathrm{Ce}^{\mathrm{M}t}\) .

\begin{enumerate}
\def\labelenumi{\alph{enumi})}
\setcounter{enumi}{4}
\item
  The map \(t \mapsto s(t)\) is a continuous semigroup on E.
\item
  The partial operator u is the infinitesimal generator of the continuous semigroup s (“Hille–Yosida theorem”). (First verify that the infinitesimal generator of s is an extension of u.)
\end{enumerate}

\(g)\) Let \(\widetilde{s}\) a continuous semigroup whose \(u\) is the infinitesimal generator. We have \(\widetilde{s} = s\) .

\begin{enumerate}
\def\labelenumi{\arabic{enumi})}
\setcounter{enumi}{6}
\tightlist
\item
  Let E be a complex Banach space and u a partial operator on E. We say that u is accretive if, for every \(x \in \text{dom}(u)\) there exists a continuous linear form \(\ell \in E'\) such that
\end{enumerate}

\begin{enumerate}
\def\labelenumi{(\roman{enumi})}
\item
  \(\mathcal{R}(\ell(u(x))) \geqslant 0;\)
\item
  One has \(\|\ell\|=\|x\|\) and \(\ell(x)=\|x\|^{2}\) .
\end{enumerate}

We say that a continuous semigroup \(s\) on E is contractive if \(\|s(t)\|\leqslant1\) for all \(t\in\mathbf{R}_{+}\) .

\begin{enumerate}
\def\labelenumi{\alph{enumi})}
\item
  The partial operator \(u\) is the infinitesimal generator of a continuous contractive semigroup on E if and only if \(u\) is accretive and if there exists \(\lambda > 0\) such that the partial operator \(\lambda 1_{\mathrm{E}} + u\) is surjective. (To show that the condition is necessary, let \(s\) be the contractive semigroup whose infinitesimal generator is \(u\) ; for \(x \in \operatorname{dom}(u)\) , consider the function defined by \(t \mapsto \ell(s(t)x)\) ; it is differentiable at 0, and its derivative at 0 is \(-\mathcal{R}(\ell(u(x)))\) .)
\item
  Suppose that u is closed and that \(u\) and \(^{t}u\) are accretive. Then u is the infinitesimal generator of a continuous contractive semigroup on E.
\end{enumerate}

\begin{enumerate}
\def\labelenumi{\arabic{enumi})}
\setcounter{enumi}{7}
\tightlist
\item
  Let E be a complex Banach space. Let s be a continuous contractive semigroup on E, and let u be its infinitesimal generator. We have
\end{enumerate}

\[
\lim _ {n \to + \infty} \left(1 + \frac {t}{n} u\right) ^ {- n} x = s (t) x
\]

for all \(x \in E\) .

\begin{enumerate}
\def\labelenumi{\arabic{enumi})}
\setcounter{enumi}{8}
\tightlist
\item
  \begin{enumerate}
  \def\labelenumii{\alph{enumii})}
  \tightlist
  \item
    Find Hilbert spaces E and F and partial operators u and v from E into F such that u and \(u + v\) are densely defined and \(u^{*} + v^{*}\) is different from \((u + v)^{*}\) .
  \end{enumerate}
\end{enumerate}

\begin{enumerate}
\def\labelenumi{\alph{enumi})}
\setcounter{enumi}{1}
\item
  Find Hilbert spaces E, F, and G and densely defined partial operators u and v from E into F and from F into G, respectively, such that \(v \circ u\) is densely defined and \(u^{*} \circ v^{*}\) is different from \((v \circ u)^{*}\) .
\item
  Find a Hilbert space E and partial self-adjoint operators u and v on E such that \(u + v\) is densely defined but is not self-adjoint.
\end{enumerate}

\begin{enumerate}
\def\labelenumi{\arabic{enumi})}
\setcounter{enumi}{9}
\item
  Give an example of multiplication operators. \(m_{g_{1}}\) and \(m_{g_{2}}\) such that \(m_{g_{1}} \circ m_{g_{2}} \neq m_{g_{1}g_{2}}\) .
\item
  We identify the space \(\mathcal{D}(]0,1[)\) with a subspace of \(\mathrm{L}^2 ([0,1])\) Let \(u\) the partial operator on \(\mathrm{L}^2 ([0,1])\) with domain \(\mathcal{D}(]0,1[)\) such that \(u(f) = i^{-1}f'\) for every function \(f \in \mathcal{D}(]0,1[)\) .
\end{enumerate}

\begin{enumerate}
\def\labelenumi{\alph{enumi})}
\item
  The partial operator \(u\) is closable and \(\overline{u}\) is symmetric.
\item
  The eigenspaces of \(\overline{u}\) relative to i and to -i are not orthogonal.
\end{enumerate}

\begin{enumerate}
\def\labelenumi{\arabic{enumi})}
\setcounter{enumi}{11}
\item
  Let \(S \subset \mathbf{C}\) a closed subset. Give an example of a closed operator with dense domain \(u\) on a Hilbert space \(E\) such that \(\operatorname{Sp}(u) = S\) .
\item
  Let \(u\) a non-closable partial operator on a Hilbert space E. There does not exist a complex number \(\lambda\) such that the linear map from \(\text{dom}(u)\) to E defined by \(x \mapsto \lambda x - u(x)\) is bijective and has a continuous inverse.
\item
  Let \(u\) a symmetric partial operator on a complex Hilbert space E. The map \(\lambda \mapsto \dim \operatorname{Ker}(\lambda 1_{\mathrm{E}} - u^{*})\) , of \(\mathbf{C}\) in \(\overline{\mathbf{N}}\) is locally constant on the resolvent set of \(u\) .
\item
  Let E be the complex Hilbert space \(\mathbf{C}^2\) and let \(u \in \mathcal{L}(\mathrm{E})\) the endomorphism such that \(u(ae_1 + be_2) = be_1\) .
\end{enumerate}

\begin{enumerate}
\def\labelenumi{\alph{enumi})}
\tightlist
\item
  Show that the spectrum of \(u\) is reduced to 0 and that, for all \(\lambda \in \mathbf{C}^*\) , we have the formula
\end{enumerate}

\[
\| \mathrm{R} (u, \lambda) \| ^ {2} = \frac {2}{1 + 2 | \lambda | ^ {2} - \sqrt {1 + 4 | \lambda | ^ {2}}}.
\]

Deduce the form of the pseudo-spectrum from this. \(\mathrm{PSp}_{\varepsilon}(u)\) for all \(\varepsilon > 0\) .

\begin{enumerate}
\def\labelenumi{\arabic{enumi})}
\setcounter{enumi}{15}
\item
  Find an example showing that the conclusion of Prop. 21 of IV, p. 251, is not valid in general for bounded components of the spectrum.
\item
  \begin{enumerate}
  \def\labelenumii{\alph{enumii})}
  \tightlist
  \item
    Let U be a connected open subset of C and let V ⊂ U be an open set. Let E be a complex Banach space. Let \(f\colon\) U → E be a holomorphic map, \(z_0 \in\) U and M \textgreater{} 0 such that \(\|f(z)\| \leqslant\) M for all \(z \in\) V and such that \(\|f(z_0)\| <\) M. We then have \(\|f(z)\| <\) M for all \(z \in\) U.
  \end{enumerate}
\end{enumerate}

\begin{enumerate}
\def\labelenumi{\alph{enumi})}
\setcounter{enumi}{1}
\tightlist
\item
  Let E be a complex Banach space and \(u \in \mathcal{L}(\mathrm{E})\) Let \(U \subset C\) an open subset of the unbounded connected component of \(\mathbf{C}-\mathrm{Sp}(u)\) If M \textgreater{} 0 is such that \(\|R(u, \lambda)\| \leqslant M\) for all \(\lambda \in U\) , then \(\|R(u, \lambda)\| < M\) for all \(\lambda \in U\) . In particular, if \(\mathbf{C}-\mathrm{Sp}(u)\) is connected, then the norm of the resolvent of u cannot be constant on a nonempty open subset of \(\mathbf{C}-\mathrm{Sp}(u)\) .
\end{enumerate}

\begin{enumerate}
\def\labelenumi{\arabic{enumi})}
\setcounter{enumi}{17}
\tightlist
\item
  Let E be the complex vector space of bounded sequences \((x_{k})_{k\in\mathbf{Z}}\) of complex numbers. For every \(x=(x_{k})\) in E, we denote
\end{enumerate}

\[
\| x \| _ {\mathrm{E}} = | x _ {0} | + \sup _ {k \in \mathbf {Z} - \{0 \}} | x _ {k} |.
\]

\begin{enumerate}
\def\labelenumi{\alph{enumi})}
\item
  The map \(x \mapsto \|x\|_{\mathrm{E}}\) is a norm on E; the space E equipped with this norm is a Banach space.
\item
  Let M \textgreater{} 2 be a real number. Let \(\lambda = (\lambda_k)_{k \in \mathbf{Z}}\) the sequence defined by \(\lambda_0 = M^{-1}\) and \(\lambda_k = 1\) if \(k \neq 0\) . We denote by u the linear map from E to E such that \(u((x_k)_{k \in \mathbf{Z}}) = (\lambda_k x_{k+1})_{k \in \mathbf{Z}}\) . This is a continuous endomorphism of E.
\item
  The open disk with center 0 and radius 1/M is contained in the resolvent set of u.
\end{enumerate}

\(d)\) For every \(\lambda\in\mathbf{C}\) such that \(|\lambda|<\inf(M^{-1},\frac{1}{2}-M^{-1})\) , we have \(\|\mathrm{R}(u,\lambda)\|\leqslant\mathrm{M}\) . (Use formula (6) of IV, p. 244).

\begin{enumerate}
\def\labelenumi{\alph{enumi})}
\setcounter{enumi}{4}
\tightlist
\item
  Let \(e_{0} \in E\) the sequence \((x_{k})_{k \in \mathbf{N}}\) defined by \(x_{0} = 1\) and \(x_{k} = 0\) if \(k \neq 0\) Let \(\lambda \in C\) such that \(|\lambda| < \inf(M^{-1}, \frac{1}{2} - M^{-1})\) . We have \(\|R(u, \lambda)e_{0}\|_{E} = M\) , hence \(\|R(u, \lambda)\| = M\) .
\end{enumerate}

(This example, like that of the following exercise, is due to E. Shargorodsky, “On the level sets of the resolvent norm of a linear operator”, Bull. London Math. Soc. 40 (2008), 493–504.)

\begin{enumerate}
\def\labelenumi{\arabic{enumi})}
\setcounter{enumi}{18}
\tightlist
\item
  Let E be the Hilbert space \(\ell_2(\mathbf{N}^*)\) Let \((\alpha_n)_{n \in \mathbf{N}^*}\) a sequence of real numbers \(\geqslant 2\) such that \(\alpha_n\) tends to \(+\infty\) as \(n \to +\infty\) . Set \(\beta_n = 1 + \alpha_n^{-1}\) for all \(n \in \mathbf{N}^*\) .
\end{enumerate}

\begin{enumerate}
\def\labelenumi{\alph{enumi})}
\tightlist
\item
  Let \(\mathrm{F}\subset \mathrm{E}\) the space of sequences \((x_{k})_{k\in \mathbf{N}^{*}}\in \mathrm{E}\) such that the series
\end{enumerate}

\[
\sum_ {n \in \mathbf {N} ^ {*}} \alpha_ {n} ^ {2} | x _ {2 n} | ^ {2}
\]

converges. It is a dense subspace of E.

\begin{enumerate}
\def\labelenumi{\alph{enumi})}
\setcounter{enumi}{1}
\tightlist
\item
  There exists a closed partial operator. \(u\) on E with domain F such that for every \(x = (x_{k})_{k \in \mathbf{N}^{*}}\) in F, we have \(u(x) = (y_{k})\) where
\end{enumerate}

\[
\binom{y _ {2 k - 1}}{y _ {2 k}} = \left( \begin{array}{c c} 0 & \alpha_ {k} \\ \beta_ {k} & 0 \end{array} \right) \binom{x _ {2 k - 1}}{x _ {2 k}}
\]

for all \(k \in \mathbf{N}^*\) .

\begin{enumerate}
\def\labelenumi{\alph{enumi})}
\setcounter{enumi}{2}
\tightlist
\item
  Let \(u_{n}\) the endomorphism
\end{enumerate}

\[
x \mapsto \left( \begin{array}{c c} 0 & \alpha_ {n} \\ \beta_ {n} & 0 \end{array} \right) x
\]

of the Hilbertian space \(C^{2}\) equipped with the canonical inner product. For every \(\lambda \in C\) such that \(|\lambda| < 1\) , the resolvent of \(u_{n}\) at \(\lambda\) is defined and we have

\[
\lim _ {n \to + \infty} \| \mathrm{R} (u _ {n}, \lambda) \| = 1.
\]

\begin{enumerate}
\def\labelenumi{\alph{enumi})}
\setcounter{enumi}{3}
\tightlist
\item
  For every \(n \in N^{*}\) and every \(\lambda \in C\) such that \(|\lambda| < \frac{1}{2}\) , we have \(\|R(u_{n}, \lambda)\| < 1\) .
\item
  The open disk with center 0 and radius \(\frac{1}{2}\) is contained in the resolvent set of u, and we have \(\|R(u,\lambda)\|=1\) for all \(\lambda\in C\) such that \(|\lambda|<\frac{1}{2}\) .
\end{enumerate}

\begin{enumerate}
\def\labelenumi{\arabic{enumi})}
\setcounter{enumi}{19}
\tightlist
\item
  If \(a\) and \(b\) are real numbers with \(a < b\) , we denote \(\mathcal{C}^2([a,b])\) the vector space of continuous functions on \([a,b]\) whose restriction to \(]a,b[\) is of class \(\mathrm{C}^2\) . We denote \(\mu\) Lebesgue measure.
\end{enumerate}

\begin{enumerate}
\def\labelenumi{\alph{enumi})}
\tightlist
\item
  Let \(q_{1}\) and \(q_{2}\) of real-valued functions in \(\mathcal{C}^{2}([a,b])\) such that \(q_{1} \geqslant q_{2}\) . For \(i \in \{1,2\}\) , let \(f_{i} \in \mathcal{C}^{2}([a,b])\) a solution of the equation
\end{enumerate}

\[
f _ {i} ^ {\prime \prime} + q _ {i} f _ {i} = 0.
\]

Suppose that \(f_{2}(a) = f_{2}(b) = 0\) and that \(f_{2} \neq 0\) If the function \(f_{1}\) does not vanish in \(]a, b[\) , then \(q_{1} = q_{2}\) and there exists \(c \in \mathbf{C}\) such that \(f_{1} = cf_{2}\) . (Reduce to the case where \(f_1\) and \(f_2\) are strictly positive on \(]a, b[;\) consider the Wronskian \(w = f_1f_2' - f_1'f_2\) and prove that \(w(a) = w(b)\) then that \(w = 0.\) )

\begin{enumerate}
\def\labelenumi{\alph{enumi})}
\setcounter{enumi}{1}
\tightlist
\item
  Let \(q \in \mathcal{C}^2([0,1])\) satisfying \(q \geqslant 0\) . We denote \(\widetilde{\mathrm{SL}}_q\) the partial differential operator on \(\mathrm{L}^2([0,1],\mu)\) with domain
\end{enumerate}

\[
\mathrm{D} = \{f \in \mathcal {C} ^ {2} ([ 0, 1 ]) \mid f (0) = f (1) = 0 \}
\]

such that

\[
\widetilde {\mathrm{SL}} _ {q} (f) = - f ^ {\prime \prime} + q f
\]

Sturm–Liouville operator. The operator \(\widetilde{SL}_{q}\) is symmetric and essentially self-adjoint. We denote \(SL_{q}\) the closure of \(\widetilde{SL}_{q}\) .

\begin{enumerate}
\def\labelenumi{\alph{enumi})}
\setcounter{enumi}{2}
\tightlist
\item
  Let \(\lambda \in \mathrm{Sp}(\mathrm{SL}_q)\) an eigenvalue of \(\mathrm{SL}_q\) The eigenspace associated with \(\lambda\) has dimension 1 and is generated by a real-valued function belonging to D.
\end{enumerate}

\(d\) ) Every eigenvalue \(\lambda\) of \(\mathrm{SL}_q\) satisfies \(\lambda \geqslant \pi^2\) . (First prove that \(\lambda\) is positive using question \(a\) ) with \(q_{1} = q\) and \(q_{2} = 0\) .)

\begin{enumerate}
\def\labelenumi{\alph{enumi})}
\setcounter{enumi}{4}
\tightlist
\item
  For s \textgreater{} 0, there exists a unique function \(f_{s} \in \mathcal{C}^{2}([0,1])\) which is a solution of the linear differential equation
\end{enumerate}

\[
- f _ {s} ^ {\prime \prime} + q f - s ^ {2} f = 0
\]

and satisfies \(f_{s}(0)=0\) and \(f_{s}^{\prime}(0)=s\) . Set \(f(x,s)=f_{s}(x)\) ; the map f is a function of class \(C^{2}\) onto {]}0,1{[}× \(R_{+}^{*}\) . The eigenvalues of \(SL_{q}\) are of the form \(\lambda=s^{2}\) where s satisfies \(f(1,s)=0\) .

\(f)\) For \(t \in [0,1]\) and \(s > 0\) , let \(\mathrm{N}(t,s) \in \overline{\mathbf{N}}\) the number of \(x \in [0,t]\) such that \(f(x,s) = 0\) . We have \(\mathrm{N}(t,s) \in \mathbf{N}\) and the function \(s \mapsto \mathrm{N}(t,s)\) is increasing for every \(t\) fixed. (Use question \(a\) ).)

\(g)\) Let \(t \in [0,1]\) and \(s_0 > 0\) Let \(m = \mathrm{N}(t,s_0) \in \mathbf{N}\) . Denote by

\[
0 = x _ {1} (s _ {0}) <   x _ {2} (s _ {0}) <   \dots <   x _ {m} (s _ {0})
\]

the zeros of \(f_{s_0}\) in \([0,t]\) . For \(s \geqslant 0\) let

\[
0 = x _ {1} (s) <   \dots <   x _ {m} (s)
\]

the \(m\) first zeros of \(f_{s}\) in \([0,t]\) . The functions \(s\mapsto x_{j}(s)\) are of class \(\mathbf{C}^1\) and we have

\[
\partial_ {1} (x _ {j} (s), s) x _ {j} ^ {\prime} (s) = - \partial_ {2} f (x _ {j} (s), s).
\]

\begin{enumerate}
\def\labelenumi{\alph{enumi})}
\setcounter{enumi}{7}
\tightlist
\item
  The set of eigenvalues of \(SL_{q}\) is the set of values of an increasing sequence of strictly positive real numbers
\end{enumerate}

\[
0 <   \lambda_ {0} <   \lambda_ {1} <   \dots <   \lambda_ {n} <   \dots
\]

tending toward \(+\infty\) . Moreover, the eigenspace corresponding to \(\lambda_{n}\) has dimension 1; if \(\varphi_{n}\) of which is a basis, we have \(\varphi_{n} \in \mathcal{C}^{2}([0,1])\) and the function \(\varphi_{n}\) vanishes at exactly \(n\) distinct points of \([0,1]\) . (To prove that there exists an infinite number of eigenvalues of \(\mathrm{SL}_{q}\) to use the question \(a\) ) with \(q_{1} = M - s^{2}\) , where \(M \in R\) is well chosen, to verify that \(N(1, s)\) tends to \(+\infty\) as \(s \to +\infty\) .

\begin{enumerate}
\def\labelenumi{\arabic{enumi})}
\setcounter{enumi}{20}
\tightlist
\item
  We resume the notation of the preceding exercise. We denote by \(\varphi_{n}\) the unique eigenfunction of \(\mathrm{SL}_{q}\) for the eigenvalue \(\lambda_{n}\) such that \(\varphi_{n}^{\prime}(0)=\lambda_{n}\) .
\end{enumerate}

\begin{enumerate}
\def\labelenumi{\alph{enumi})}
\tightlist
\item
  Let \(\mathrm{M} = \sup \{q(x)\}\) . For every \(n \in \mathbf{N}\) , we have
\end{enumerate}

\[
\pi^ {2} (n + 1) ^ {2} \leqslant \lambda_ {n} \leqslant \pi^ {2} (n + 1) ^ {2} + \mathrm{M},
\]

and moreover

\[
\lambda_ {n} = \pi^ {2} (n + 1) ^ {2} \Bigl (1 + \mathrm{O} \Bigl (\frac {1}{n ^ {2}} \Bigr) \Bigr)
\]

as \(n \rightarrow +\infty\) .

\begin{enumerate}
\def\labelenumi{\alph{enumi})}
\setcounter{enumi}{1}
\tightlist
\item
  Let \(s > 0\) Let \(f \in \mathcal{C}^2([0,1])\) such that
\end{enumerate}

\[
- f ^ {\prime \prime} + s ^ {2} f = - q f, \qquad f (0) = 0, \quad f ^ {\prime} (0) = s.
\]

We then have \((1-s^{-1}u_{s})f=g_{s}\) , where \(g_{s}(x)=\sin(sx)\) and \(u_{s}\) is the integral operator on \(\mathcal{C}([0,1])\) given by

\[
u _ {s} (f) (x) = \int_ {0} ^ {x} \sin (s (x - y)) q (y) f (y) d \mu (y).
\]

\begin{enumerate}
\def\labelenumi{\alph{enumi})}
\setcounter{enumi}{2}
\item
  For \(s\) sufficiently large, we have \(s \notin \operatorname{Sp}(u_s)\) .
\item
  For every fixed x, we have
\end{enumerate}

\[
\varphi_ {n} (x) = \sin ((n + 1) x) + \mathrm{O} (n ^ {- 1})
\]

as \(n \rightarrow +\infty\) .

\begin{enumerate}
\def\labelenumi{\arabic{enumi})}
\setcounter{enumi}{21}
\tightlist
\item
  We resume the notation of the preceding exercise.
\end{enumerate}

\begin{enumerate}
\def\labelenumi{\alph{enumi})}
\tightlist
\item
  There exist functions \(f_{1}\) and \(f_{2}\) in \(\mathcal{C}^{2}([a,b])\) such that \(-f_{i}^{\prime\prime} + qf_{i} = 0\) and
\end{enumerate}

\[
\begin{array}{c c} f _ {1} (0) = 0, & f _ {1} ^ {\prime} (0) = 1 \\ f _ {2} (1) = 0, & f _ {2} ^ {\prime} (1) = - 1. \end{array}
\]

These functions are the unique solutions of these equations, and the function \(w = f_{1}f_{2}' - f_{1}'f_{2}\) is constant.

\begin{enumerate}
\def\labelenumi{\alph{enumi})}
\setcounter{enumi}{1}
\tightlist
\item
  Let \(g \in \mathcal{C}([0,1])\) . There exists a unique function \(\varphi \in \mathcal{C}^2([0,1])\) such that
\end{enumerate}

\[
- \varphi^ {\prime \prime} + q \varphi = g,
\]

and we have \(\varphi = c_{1}f_{1} + c_{2}f_{2}\) , where

\[
c _ {1} (x) = - \frac {1}{w} \int_ {x} ^ {1} g (y) f _ {2} (y) d \mu (y), \qquad c _ {2} (x) = - \frac {1}{w} \int_ {0} ^ {x} g (y) f _ {1} (y) d \mu (y)
\]

for all \(x \in [0,1]\) .

\begin{enumerate}
\def\labelenumi{\alph{enumi})}
\setcounter{enumi}{2}
\tightlist
\item
  The partial operator \(\mathrm{SL}_q\) is surjective; its inverse is the map \(g\mapsto \varphi\) and there exists a function \(\mathrm{G}\in \mathcal{C}([0,1]^2)\) such that
\end{enumerate}

\[
\mathrm{SL} _ {q} ^ {- 1} (g) (x) = \int_ {0} ^ {1} \mathrm{G} (x, y) g (y) d \mu (y)
\]

for all \(g \in \mathcal{C}([0,1])\) and almost all \(x \in [0,1]\) . We have \(G(x,y) = G(y,x)\) for all \((x,y) \in [0,1]^2\) .

\(d)\) The operator \(\mathrm{SL}_q^{-1}\) is compact and injective on \(\mathrm{L}^2 ([0,1],\mu)\) .

\begin{enumerate}
\def\labelenumi{\alph{enumi})}
\setcounter{enumi}{4}
\tightlist
\item
  The functions \((\varphi_{n})\) form an orthonormal basis of \(\mathrm{L}^{2}([0,1])\) , and the spectrum of \(SL_{q}\) coincides with the set of its eigenvalues.
\end{enumerate}

\begin{enumerate}
\def\labelenumi{\arabic{enumi})}
\setcounter{enumi}{22}
\tightlist
\item
  Let \(E = \ell_{\mathbf{C}}^{2}(\mathbf{N})\) and let u be the endomorphism of E defined by
\end{enumerate}

\[
u ((x _ {i}) _ {i \in \mathbf {N}}) = (0, x _ {1}, \dots , x _ {n}, \dots)
\]

for all \((x_{i})\in \mathrm{E}\)

\begin{enumerate}
\def\labelenumi{\alph{enumi})}
\item
  The spectrum of \(u\) is the closed unit disk in \(\mathbf{C}\) .
\item
  The residual spectrum of u is the open unit disk in C.
\item
  The essential spectrum of \(u\) is the unit circle in \(\mathbf{C}\) .
\end{enumerate}

\begin{enumerate}
\def\labelenumi{\arabic{enumi})}
\setcounter{enumi}{23}
\tightlist
\item
  Let E be a complex Hilbert space and u a partial symmetric operator on E. Let \(F = E \oplus E\) and let v be the partial operator on F with domain \(\text{dom}(u^{*}) \oplus \text{dom}(\overline{u})\) such that \(v(x, y) = (\overline{u}(y), u^{*}(x))\) .
\end{enumerate}

\begin{enumerate}
\def\labelenumi{\alph{enumi})}
\item
  The partial operator v is self-adjoint.
\item
  Let \(\widetilde{\mathrm{E}}\) the image of the linear map \(\delta\colon x\mapsto(x,x)\) The reduction of \(v\) to \(\widetilde{\mathrm{E}}\) is the partial operator with domain \(\delta(\mathrm{dom}(u))\) such that \((x,x)\mapsto(u(x),u(x))\) .
\end{enumerate}

\exercisesection{}

Unless stated otherwise, the Hilbert spaces below are assumed to be complex.

\begin{enumerate}
\def\labelenumi{\arabic{enumi})}
\tightlist
\item
  Let E be a Hilbert space and \(\mathcal{P}(\mathrm{E})\subset \mathcal{L}(\mathrm{E})\) the set of orthogonal projections in E. Let S be a universally measurable subset of C. We denote by \(\mathcal{T}(\mathrm{S})\) the set of universally measurable subsets T of C such that \(\mathrm{T}\subset \mathrm{S}\) . A measure on S with values in the projections of E is a map \(\varpi\) of \(\mathcal{T}(\mathrm{S})\) in \(\mathcal{P}(\mathrm{E})\) such that
\end{enumerate}

\begin{enumerate}
\def\labelenumi{(\roman{enumi})}
\item
  One has \(\varpi(\varnothing)=0\) and \(\varpi(S)=1_{E}\) .
\item
  If \((\mathrm{T}_{n})\) is a sequence of pairwise disjoint sets in \(\mathcal{T}(\mathrm{S})\) and if T is the union of the sets \(T_{n}\) , then
\end{enumerate}

\[
\varpi (\mathrm{T}) = \sum_ {n \in \mathbf {N}} \varpi (\mathrm{T} _ {n})
\]

in the space \(\mathcal{L}(\mathrm{E})\) endowed with the topology of simple convergence.

For every normal partial operator u on E with spectrum equal to S, and for \(T \in \mathcal{T}(S)\) , we denote \(p_{T,u}\) the spectral projector of u defined by T.

The map which to a normal partial operator \(u\) on E with spectrum S associates the map \(\mathrm{T} \mapsto p_{\mathrm{T}, u}\) is a bijection between the set of normal partial operators with spectrum S and the set of measures on S with values in the projections of E.

\begin{enumerate}
\def\labelenumi{\arabic{enumi})}
\setcounter{enumi}{1}
\tightlist
\item
  Let E be a complex Hilbert space and let u be a normal partial operator on E. We denote by \(p_{A}\) the spectral projection of u on A. Let a \textless{} b be real numbers. We have
\end{enumerate}

\[
\frac {1}{2} (p _ {[ a, b ]} + p _ {] a, b [}) = \lim _ {\varepsilon \rightarrow 0} \frac {1}{2 i \pi} \int_ {a} ^ {b} \big (\mathrm{R} (u, t - i \varepsilon) - \mathrm{R} (u, t + i \varepsilon) \big) d t,
\]

in the space \(\mathcal{L}(\mathrm{E})\) endowed with the topology of simple convergence (“Stone formula”).

\begin{enumerate}
\def\labelenumi{\arabic{enumi})}
\setcounter{enumi}{2}
\item
  Let E be a Hilbert space of countable type. The only closed two-sided ideals of \(\mathcal{L}(\mathrm{E})\) are \(\{0\}\) , \(\mathcal{L}^{\mathrm{c}}(\mathrm{E})\) and \(\mathcal{L}(\mathrm{E})\) . (Use exerc. IV, p. 316, exerc. 10; if I is such an ideal, suppose that there exists in I a nonzero self-adjoint element \(u\) which is not compact, and show that for every open interval J in R such that \(0 \notin J\) , the spectral projection of \(u\) on J belongs to I.)
\item
  Let E be a Hilbert space and u a self-adjoint partial operator on E. The map \(\sigma\colon f\mapsto f(u)\) of \(\mathcal{L}_{\mathrm{u}}^{\infty}(\mathrm{Sp}(u))\) in \(\mathcal{L}(\mathrm{E})\) is the unique morphism of involutive algebras satisfying the following three conditions:
\end{enumerate}

\begin{enumerate}
\def\labelenumi{(\roman{enumi})}
\item
  One has \(\| \sigma (f)\| \leqslant \| f\|_{\infty}\) for all \(f\in \mathcal{L}_{\mathrm{u}}^{\infty}(\mathrm{Sp}(u))\)
\item
  If \((f_n)_{n\in \mathbf{N}}\) is a bounded sequence in \(\mathcal{L}_{\mathrm{u}}^{\infty}(\mathrm{Sp}(u))\) and \(f_{n}(t)\to f(t)\) for all \(t\in \mathrm{Sp}(u)\) , then \(\sigma (f_n)\rightarrow \sigma (f)\) in \(\mathcal{L}(\mathrm{E})\) endowed with the topology of simple convergence.
\item
  If \((f_n)_{n\in \mathbf{N}}\) is a sequence in \(\mathcal{L}_{\mathrm{u}}^{\infty}(\mathrm{Sp}(u))\) such that \(|f_n(t)|\leqslant |t|\) and \(f_{n}(t)\to t\) for all \(t\in \mathrm{Sp}(u)\) , then we have \(u(x) = \lim \sigma (f_n)(x)\) for all \(x\in \mathrm{dom}(u)\) .
\end{enumerate}

\begin{enumerate}
\def\labelenumi{\arabic{enumi})}
\setcounter{enumi}{4}
\item
  Let u be a normal partial operator on a Hilbert space E. The residual spectrum of u is empty.
\item
  Let E and F be Hilbert spaces. Let u be a closed partial operator with dense domain from E to F. Let \(\lambda \in C^{*}\) belonging to the resolvent set of the self-adjoint operator \(u^{*} \circ u\) on E (cf. prop. 12 of IV, p. 241).
\end{enumerate}

\begin{enumerate}
\def\labelenumi{\alph{enumi})}
\item
  The partial operator \(u \circ u^{*}\) from F to F is self-adjoint.
\item
  For every vector \(x \in \mathrm{dom}(u^{*})\) , the vector
\end{enumerate}

\[
y = (1 _ {\mathrm{F}} + u \circ \mathrm{R} (u ^ {*} \circ u, \lambda) \circ u ^ {*}) x
\]

belongs to dom \((u\circ u^{*})\) , and the formula

\[
(u \circ u ^ {*} - \lambda 1 _ {\mathrm{F}}) y = x
\]

is valid.

\begin{enumerate}
\def\labelenumi{\alph{enumi})}
\setcounter{enumi}{2}
\tightlist
\item
  The partial operator \(u \circ \mathrm{R}(u^* \circ u, \lambda) \circ u^*\) has dense domain, and is continuous on its domain.
\end{enumerate}

\(d)\) The partial operator \(u \circ u^{*} - \lambda 1_{\mathrm{F}}\) is surjective.

\begin{enumerate}
\def\labelenumi{\alph{enumi})}
\setcounter{enumi}{4}
\item
  One has \(\lambda \notin \operatorname{Sp}(u \circ u^{*})\) .
\item
  We have the identity
\end{enumerate}

\[
\lambda \mathrm{R} (u \circ u ^ {*}, \lambda) - u \circ \mathrm{R} (u ^ {*} \circ u, \lambda) \circ u ^ {*} = 1 _ {\mathrm{F}}
\]

in \(\mathcal{L}(\mathrm{F})\) (where \(u\circ \mathrm{R}(u^{*}\circ u,\lambda)\circ u^{*}\) denotes, by abuse of notation, the endomorphism of F which coincides with this partial operator on its domain, cf. question \(c\) )).

\begin{enumerate}
\def\labelenumi{\alph{enumi})}
\setcounter{enumi}{6}
\tightlist
\item
  We have the relation
\end{enumerate}

\[
u ^ {*} \circ \mathrm{R} (u \circ u ^ {*}, \lambda) \subset \mathrm{R} (u ^ {*} \circ u, \lambda) \circ u ^ {*}.
\]

\begin{enumerate}
\def\labelenumi{\alph{enumi})}
\setcounter{enumi}{7}
\tightlist
\item
  The partial operator u induces, by passage to subspaces, an isomorphism of \(\operatorname{Ker}(u^{*} \circ u - \lambda1_{\mathrm{E}})\) in \(\operatorname{Ker}(u \circ u^{*} - \lambda1_{\mathrm{F}})\) .
\end{enumerate}

\begin{enumerate}
\def\labelenumi{\arabic{enumi})}
\setcounter{enumi}{6}
\tightlist
\item
  Let A be a set of closed operators with dense domain on a Hilbert space E. Let \(\lambda \in \mathbf{C}\) . We say that A is \(\lambda\) -commutative if \(\lambda\) belongs to the resolvent set of \(u\) for all \(u \in \mathrm{A}\) , and if the endomorphisms \(\mathrm{R}(u, \lambda)\) of E, for \(u \in \mathrm{A}\) , are pairwise permutable.
\end{enumerate}

We say that A is commutative in the resolvent sense if there exists \(\lambda \in \mathbf{C}\) such that A is \(\lambda\) -commutative.

A set consisting of a single closed densely defined operator u is commutative if and only if the resolvent set of u is nonempty.

\begin{enumerate}
\def\labelenumi{\alph{enumi})}
\item
  If A ⊂ L(E), then A is commutative in the resolvent sense if and only if the union of the spectra of the elements of A is not equal to C, and A is commutative.
\item
  Let A be a set of closed densely defined operators on E that is λ-commutative. Then A is μ-commutative for every complex number μ belonging to the resolvent set of every \(u \in A\) .
\item
  Let A be a set of closed densely defined operators on E, commutative in the resolvent sense, and such that \(u^{*} \in A\) for all \(u \in A\) There exists a locally compact topological space X, a positive measure \(\mu\) on X, an isometric isomorphism \(v\) of \(\mathrm{L}^2 (\mathrm{X},\mu)\) on E, and a mapping \(u\mapsto g_u\) from A into the set of functions \(\mu\) -measurable such that \(u = v\circ m_{g_u}\circ v^{-1}\) for all \(u\in \mathrm{A}\) .
\end{enumerate}

\begin{enumerate}
\def\labelenumi{\arabic{enumi})}
\setcounter{enumi}{7}
\tightlist
\item
  Let u and v be partial normal operators on a Hilbert space E. The following assertions are equivalent:
\end{enumerate}

\begin{enumerate}
\def\labelenumi{\alph{enumi})}
\item
  The endomorphisms \(b(u)\) and \(b(v)\) are permutable;
\item
  For every universally measurable subset A of C, the spectral projections of u and v on A commute;
\end{enumerate}

If the resolvent set of \(u\) and \(v\) is nonempty, these assertions are equivalent to

\begin{enumerate}
\def\labelenumi{\alph{enumi})}
\setcounter{enumi}{2}
\tightlist
\item
  For every \(\lambda \notin \operatorname{Sp}(u)\) and \(\lambda' \notin \operatorname{Sp}(v)\) , the resolvents \(\operatorname{R}(u, \lambda)\) and \(\operatorname{R}(v, \lambda')\) commute.
\end{enumerate}

If \(u\) and \(v\) are self-adjoint, these assertions are equivalent to

\begin{enumerate}
\def\labelenumi{\alph{enumi})}
\setcounter{enumi}{3}
\tightlist
\item
  For every \(t \in R\) , the unitary endomorphisms \(e^{itu}\) and \(e^{itv}\) are permutable.
\end{enumerate}

\begin{enumerate}
\def\labelenumi{\arabic{enumi})}
\setcounter{enumi}{8}
\tightlist
\item
  Let u be a closed densely defined partial operator on a Hilbert space E. The following assertions are equivalent:
\end{enumerate}

\begin{enumerate}
\def\labelenumi{\alph{enumi})}
\item
  The operator u is normal;
\item
  There exist self-adjoint operators \(u_{1}\) and \(u_{2}\) on E, which commute in the resolvent sense, and such that \(u = u_{1} + iu_{2}\) ;
\item
  One has \(\operatorname{dom}(u) = \operatorname{dom}(u^*)\) and \(\| u(x)\| = \| u^{*}(x)\|\) for all \(x\in \operatorname{dom}(u)\) ;
\item
  One has \(uu^{*} = u^{*}u\) .
\end{enumerate}

If the resolvent set of \(u\) is nonempty, these assertions are equivalent to

\begin{enumerate}
\def\labelenumi{\alph{enumi})}
\setcounter{enumi}{4}
\tightlist
\item
  For every \(\lambda \notin \operatorname{Sp}(u)\) , the resolvent \(\operatorname{R}(u, \lambda)\) is normal.
\end{enumerate}

\begin{enumerate}
\def\labelenumi{\arabic{enumi})}
\setcounter{enumi}{9}
\tightlist
\item
  Let E be a complex Hilbert space and A a unital self-adjoint subalgebra of \(\mathcal{L}(\mathrm{E})\) which is closed under the topology of simple convergence. (Such an algebra is called a von Neumann algebra.)
\end{enumerate}

\begin{enumerate}
\def\labelenumi{\alph{enumi})}
\item
  There exists a subset X of \(\mathcal{L}(\mathrm{E})\) such that A coincides with the commutant of X in \(\mathcal{L}(\mathrm{E})\) . (Use Exercise 24 of IV, p. 326.) Conversely, for every self-adjoint subset of X, the commutant of X in \(\mathcal{L}(\mathrm{E})\) is a unital self-adjoint subalgebra closed under the topology of pointwise convergence.
\item
  Let \(u \in \mathrm{A}\) and let \((j, |u|)\) the polar decomposition of \(u\) (definition 4 of I, p. 140). The elements \(j\) and \(u\) belong to A.
\item
  Let \(u \in A\) such that u is normal. For every function \(f \in \mathcal{L}_{\mathrm{u}}^{\infty}(\mathrm{Sp}(u))\) , the endomorphism \(f(u)\) belongs to A.
\item
  Let \(u \in A\) . The element u is a right (resp. left) zero divisor in A (cf. A, I, p. 93) if and only if the image of u is not dense in E (resp. if u is not injective).
\item
  Let F be the set of partial operators on E whose elements are the closed partial operators with dense domain u such that, for every x in X, the partial operator \(x \circ u\) is a reduction of the partial operator \(u \circ x\) For every u in F and v in A, the partial operators \(u^{*}\) , \(u + v\) , \(v \circ u\) and \(u \circ v\) belong to F.
\end{enumerate}

\(f\) ) Let \(u\in \mathrm{F}\) and let \((j,|u|)\) the polar decomposition of \(u\) (definition 7 of IV, p. 290). The partial operator \(|u|\) belongs to F and \(j\) belongs to A.

\(g)\) Let \(u \in \mathrm{F}\) such that \(u\) is normal. For every function \(f \in \mathcal{L}_{\mathrm{u}}(\mathrm{Sp}(u))\) , the partial operator \(f(u)\) belongs to F; if \(f\) is bounded, then \(f(u)\) belongs to A.

\begin{enumerate}
\def\labelenumi{\alph{enumi})}
\setcounter{enumi}{7}
\item
  Let \(u \in A\) injective whose image is dense in E. The partial operator \(u^{-1}\) belongs to F.
\item
  For every u in F, there exist elements \(u_{1}\) and \(u_{2}\) of A such that \(u_{2}\) is injective with dense image and \(u = u_{1}u_{2}^{-1}\) . (First reduce to the case where u is self-adjoint and positive, then apply the universally measurable functional calculus to the functions \(f_{1}: t \mapsto t\varphi(t) + \psi(t)\) and \(f_{2}: t \mapsto \varphi(t) + t^{-1}\psi(t)\) , where \(\varphi\) and \(\psi\) are the characteristic functions of the intervals {[}0, 1{[} and {[}1, +∞{[} respectively.{]})
\item
  Assume that every injective element of A with dense image is surjective. The algebra A then satisfies the Ore condition: for any u and s in A, if s is not a zero divisor (neither on the right nor on the left), then there exists \(u'\) and \(s'\) in A such that \(s'\) is not a zero divisor and \(us' = su'\) .
\end{enumerate}

\(k\) ) The conclusion of the previous question is not always valid. (Consider \(\mathrm{A} = \mathcal{L}(\mathrm{E}).\) )

\begin{enumerate}
\def\labelenumi{\arabic{enumi})}
\setcounter{enumi}{10}
\tightlist
\item
  \begin{enumerate}
  \def\labelenumii{\alph{enumii})}
  \tightlist
  \item
    Let \(q\) a positive partial form on a complex Hilbert space E. Then \(q\) is closed if and only if for every sequence \((x_{n})\) of elements of \(\text{dom}(q)\) which converges to \(x\) in E and satisfies
  \end{enumerate}
\end{enumerate}

\[
\lim _ {n, m \to + \infty} q (x _ {n} - x _ {m}, x _ {n} - x _ {m}) = 0,
\]

we have \(x \in \text{dom}(q)\) and \(q(x_{n} - x, x_{n} - x) \to 0\) .

\begin{enumerate}
\def\labelenumi{\alph{enumi})}
\setcounter{enumi}{1}
\tightlist
\item
  Let \(E = L^{2}(\mathbf{R}, \mu)\) , where \(\mu\) is the Lebesgue measure. The positive Hermitian form q with domain \(\mathcal{D}(\mathbf{R})\) defined by \(q(\varphi_{1}, \varphi_{2}) = \overline{\varphi_{1}(0)}\varphi_{2}(0)\) is not closed, and there does not exist a closed partial positive form extending q.
\end{enumerate}

\begin{enumerate}
\def\labelenumi{\arabic{enumi})}
\setcounter{enumi}{11}
\tightlist
\item
  We equip \(\mathbf{R}\) of Lebesgue measure. Let E be the closed subspace of \(\mathrm{L}^2 (\mathbf{R})\) consisting of the \(f\in \mathrm{L}^2 (\mathbf{R})\) such that \(f(-x) = f(x)\) for almost every \(x\in \mathbf{R}\) Let \(\mathcal{D}_{+}(\mathbf{R})\) the space of test functions on \(\mathbf{R}\) such that \(f(x) = f(-x)\) for all \(x\in \mathbf{R}\) . We identify \(\mathcal{D}_{+}(\mathbf{R})\) with a subspace of E. We denote by \(\mathcal{F}\in \mathcal{L}(\mathrm{E})\) the endomorphism of E deduced from the Fourier transform of \(\mathrm{L}^2 (\mathbf{R})\) by passage to the subspaces.
\end{enumerate}

\begin{enumerate}
\def\labelenumi{\alph{enumi})}
\item
  For every \(f \in \mathcal{D}_{+}(\mathbf{R})\) , we have \(\mathcal{F}(f) + \mathcal{F}^{*}(f) \in \mathcal{D}_{+}(\mathbf{R})\) .
\item
  Let u be the partial operator on E whose domain is \(\mathcal{D}_{+}(\mathbf{R})\) and which is defined by \(u(f) = \mathcal{F}(f) + \mathcal{F}^{*}(f)\) for \(f \in \mathcal{D}_{+}(\mathbf{R})\) The operator u is essentially self-adjoint.
\item
  One has \(\operatorname{dom}(u^{2})=\{0\}\) (Use Exercise 28 of II, p. 273.)
\end{enumerate}

\begin{enumerate}
\def\labelenumi{\arabic{enumi})}
\setcounter{enumi}{12}
\tightlist
\item
  Let E be a complex Hilbert space and let u be a partial symmetric operator on E. Suppose there exists a real number c \textgreater{} 0 such that \(\langle x \mid u(x) \rangle \geqslant c \|x\|^{2}\) for all \(x \in \text{dom}(u)\) .
\end{enumerate}

Show that the following conditions are equivalent:

\begin{enumerate}
\def\labelenumi{(\roman{enumi})}
\item
  The partial operator \(u\) is essentially self-adjoint;
\item
  The image of u is dense in E;
\item
  The kernel of \(u^{*}\) is reduced to 0.
\end{enumerate}

\begin{enumerate}
\def\labelenumi{\arabic{enumi})}
\setcounter{enumi}{13}
\tightlist
\item
  Let \(n \in \{1, 2, 3\}\) . We denote \(\Delta\) the Laplace operator on \(\mathrm{L}^2(\mathbf{R}^n)\) , and D its domain. It is a positive partial operator.
\end{enumerate}

\begin{enumerate}
\def\labelenumi{\alph{enumi})}
\tightlist
\item
  There exists a real number \(c \geqslant 0\) such that for every \(f \in D\) , we have
\end{enumerate}

\[
\| f \| _ {\infty} \leqslant c (\| \Delta f \| + \| f \|).
\]

\begin{enumerate}
\def\labelenumi{\alph{enumi})}
\setcounter{enumi}{1}
\tightlist
\item
  For every real number \(a > 0\) , there exists a real number \(b \geqslant 0\) such that for every \(f \in \mathrm{D}\) , we have
\end{enumerate}

\[
\| f \| _ {\infty} \leqslant a \| \Delta f \| + b \| f \|.
\]

(Apply the previous question to \(x \mapsto f(\alpha x)\) , where \(\alpha \in \mathbf{R}_{+}^{*}\) .)

\begin{enumerate}
\def\labelenumi{\alph{enumi})}
\setcounter{enumi}{2}
\item
  Let \(g \in \mathcal{L}^{\infty}(\mathbf{R}^{n}) + \mathcal{L}^{2}(\mathbf{R}^{n})\) . If \(g\) is real-valued, then the partial operator \(\Delta + m_{g}\) is self-adjoint.
\item
  Let \(\alpha\) be a real number such that \(0 \leqslant \alpha < n/2\) . Let N be the Euclidean norm on \(\mathbf{R}^n\) and \(g\) the function on \(\mathbf{R}^n\) defined by \(g(x) = 1/\mathrm{N}(x)^\alpha\) if \(x \neq 0\) , and \(g(0) = 0\) . The partial operator \(\Delta + m_g\) is self-adjoint.
\end{enumerate}

\begin{enumerate}
\def\labelenumi{\arabic{enumi})}
\setcounter{enumi}{14}
\tightlist
\item
  Let \(n \in \mathbf{N}\) and let \(\Delta\) the unique Laplacian on \(\mathrm{L}^2(\mathbf{R}^n)\) . Denote by B the unit ball in \(\mathbf{R}^n\) .
\end{enumerate}

If \(g \in \mathcal{L}_{\mathrm{u}}(\mathbf{R}^{n})\) is such that the domain of \(\Delta\) is contained in the domain of \(m_{g}\) , then

\[
\sup _ {y \in \mathbf {R} ^ {n}} \int_ {\mathrm{B}} | g (x + y) | ^ {2} d x
\]

is finite. (First verify that \(m_{g}\) is \(\Delta\) -bounded.)

\begin{enumerate}
\def\labelenumi{\arabic{enumi})}
\setcounter{enumi}{15}
\tightlist
\item
  \begin{enumerate}
  \def\labelenumii{\alph{enumii})}
  \tightlist
  \item
    Let I = {]}0, 1{[}. The space H \(^{1}\) (I) is contained in C(Ī).
  \end{enumerate}
\end{enumerate}

\begin{enumerate}
\def\labelenumi{\alph{enumi})}
\setcounter{enumi}{1}
\item
  The differential operator \(u\colon f\mapsto -if'\) onto \(\mathrm{L}^2 (\mathrm{I})\) with domain \(\mathcal{D}(\mathrm{I})\) is symmetric and closable. For all \(f\in \mathrm{dom}(u^{*})\) the derived distribution \(f'\) of \(f\) belongs to \(\mathrm{L}^2 (\mathrm{I})\) and we have \(u^{*}(f) = if'\) .
\item
  The deficiency spaces of \(\overline{u}\) are of dimension 1; deduce from this an explicit description of all self-adjoint extensions of \(\overline{u}\) .
\end{enumerate}

17) Classify the Laplacians on I ={]}0, 1{[}.

¶ 18) Let n be an integer ≥ 1. We denote by D the partial operator on L²(T\^{}n) with domain C∞(T\^{}n) such that

\[
\mathrm{D} (\varphi) = - \sum_ {i = 1} ^ {n} \partial_ {i} ^ {2} \varphi
\]

for all \(\varphi\in\mathcal{C}^{\infty}(\mathbf{T}^{n})\) ,

\begin{enumerate}
\def\labelenumi{\alph{enumi})}
\item
  The partial operator D is essentially self-adjoint. We denote by \(\Delta\) its closure.
\item
  The partial operator \(\Delta\) is self-adjoint and positive.
\item
  The partial operator \(\Delta\) has compact resolvent.
\end{enumerate}

\(d)\) Let \(\lambda\in\mathbf{R}\) . We have \(\lambda\in\mathrm{Sp}(\Delta)\) if and only if there exist natural integers \(k_{1},\ldots,k_{n}\) such that \(\lambda=4\pi^{2}(k_{1}^{2}+\cdots+k_{n}^{2})\) .

\begin{enumerate}
\def\labelenumi{\alph{enumi})}
\setcounter{enumi}{4}
\tightlist
\item
  For every eigenvalue \(\lambda\) of \(\Delta\) , let \(N(\lambda)\) the dimension of the corresponding eigenspace. We have
\end{enumerate}

\[
\sum_ {\lambda \leqslant \mathrm{T}} \mathrm{N} (\lambda) = c _ {n} \mathrm{T} ^ {n / 2} + \mathrm{O} (\mathrm{T} ^ {(n - 1) / 2})
\]

when T → +∞.

\(f)\) If \(n \geqslant 2\) , then \(\limsup_{\lambda} N(\lambda) = +\infty\) . (In the case \(n = 2\) , consider \(\lambda = 4\pi^2 m\) , where \(m\) is a product of prime numbers of the form \(4k + 1\) .) \(g)\) If \(n = 1\) or \(n = 2\) , we have

\[
\operatorname{Card} (\operatorname{Sp} (u) \cap [ 0, \mathrm{T} ]) = o (\mathrm{T}),
\]

and if \(n\geqslant 3\) , we have

\[
\liminf _ {T \to + \infty} \frac {1}{T} \operatorname{Card} (\operatorname{Sp} (u) \cap [ 0, T ]) > 0.
\]

\begin{enumerate}
\def\labelenumi{\arabic{enumi})}
\setcounter{enumi}{18}
\tightlist
\item
  Let u be a partial self-adjoint operator on a Hilbert space E.
\end{enumerate}

\begin{enumerate}
\def\labelenumi{\alph{enumi})}
\tightlist
\item
  Let \(t \in R\) and \(x \in E\) . If the spectral measure of x with respect to u has compact support, then we have
\end{enumerate}

\[
e ^ {i u} (x) = \sum_ {n \in \mathbf {N}} \frac {u ^ {n} (x)}{n !},
\]

where the series converges in E.

\begin{enumerate}
\def\labelenumi{\alph{enumi})}
\setcounter{enumi}{1}
\tightlist
\item
  The set of vectors \(x \in E\) such that the spectral measure of x with respect to u has compact support is a dense subspace of E.
\end{enumerate}

\begin{enumerate}
\def\labelenumi{\arabic{enumi})}
\setcounter{enumi}{19}
\tightlist
\item
  Let \(u\) a partial operator on a Hilbert space E. We denote by \(\mathrm{dom}(u)^{\mathrm{an}}\) the space of vectors \(x \in \mathrm{E}\) such that \(x \in \mathrm{dom}(u^n)\) for every integer \(n \in \mathbf{N}\) , and such that the power series
\end{enumerate}

\[
\sum_ {n \in \mathbf {N}} \frac {\| u ^ {n} (x) \|}{n !} z ^ {n}\tag{23}
\]

has a radius of convergence \textgreater{} 0. The elements of \(\text{dom}(u)^{\text{an}}\) are said to be the analytic vectors of u.

Let \(x\) an element of E such that \(x\) belongs to \(\mathrm{dom}(u^n)\) for all \(n \in \mathbf{N}\) Let \(\mathrm{E}_x\) the closure of the subspace of E generated by the vectors \(u^n(x)\) for \(n \in \mathbf{N}\) and let \(u_x\) the partial operator on \(\overline{\mathrm{E}}_x\) with domain \(\mathrm{E}_x\) deduced from \(u\) by passing to subspaces. We say that \(x\) is a uniqueness vector of \(u\) if \(u_x\) is essentially self-adjoint.

\begin{enumerate}
\def\labelenumi{\alph{enumi})}
\tightlist
\item
  Suppose that u is symmetric and that E contains a total subset of uniqueness vectors for u. Then u is essentially self-adjoint. (Show that the image of \(u + i1_{E}\) and that of \(u - i1_{E}\) are dense in E.)
\end{enumerate}

Suppose henceforth that u is symmetric and that E contains a total subset of analytic vectors.

\begin{enumerate}
\def\labelenumi{\alph{enumi})}
\setcounter{enumi}{1}
\item
  For every analytic vector x of u, the partial operator \(u_{x}\) onto \(\overline{E}_{x}\) admits a self-adjoint extension.
\item
  Let \(x\) an analytic vector for \(u\) and \(r > 0\) the radius of convergence of the power series (23). Let \(v\) be a self-adjoint extension of \(u_x\) . For every \(t \in \mathbf{R}\) such that \(|t| < r\) , we have
\end{enumerate}

\[
\langle e ^ {i t v} (x) \mid x \rangle = \sum_ {n \in \mathbf {N}} \frac {(i t) ^ {n}}{n !} \langle u ^ {n} (x) \mid x \rangle .
\]

(Using the spectral measure of \(x\) relative to \(v\) , show that the left-hand side is the restriction of an analytic function on the set of \(t \in \mathbf{C}\) such that \(|t| < r\) , and compute the power series expansion of this function at 0.)

\begin{enumerate}
\def\labelenumi{\alph{enumi})}
\setcounter{enumi}{3}
\item
  For every \(t \in R\) , the unitary operator \(e^{itv}\) depends only on \(u_x\) and the operator \(u_x\) is essentially self-adjoint. (Use Theorem V, p. 428.) e) Deduce that u is essentially self-adjoint.
\item
  A closed symmetric partial operator u on E is self-adjoint if and only if \(\text{dom}(u)^{\text{an}}\) is dense in E.
\end{enumerate}

\begin{enumerate}
\def\labelenumi{\arabic{enumi})}
\setcounter{enumi}{20}
\tightlist
\item
  We denote \(x\) the identity function of \(\mathbf{R}\) . We denote \(\mathcal{M}_{\mathrm{m}}(\mathbf{R})\) the set of positive measures \(\nu\) onto \(\mathbf{R}\) such that \(x^n \in \mathcal{L}^1(\mathbf{R}, \nu)\) for every integer \(n \in \mathbf{N}\) . \(a)\) Let \(\nu \in \mathcal{M}_{\mathrm{m}}(\mathbf{R})\) and set \(\mu_n = \nu(x^n)\) for all \(n \in \mathbf{N}\) (“moments of the measure \(\nu\) ”). For every integer \(n \geqslant 0\) and for all complex numbers \(a_0, \ldots, a_n\) , we have
\end{enumerate}

\[
\sum_ {i = 0} ^ {n} \sum_ {j = 0} ^ {n} \overline {{a}} _ {i} a _ {j} \mu_ {i + j} \geqslant 0.\tag{24}
\]

In what follows, we propose to study the converse assertion of question \(a\) ) (“Hamburger moment problem”). In the rest of this exercise, we therefore assume given a sequence \(\boldsymbol{\mu} = (\mu_n)_{n \geqslant 0}\) of real numbers satisfying condition (24).

\begin{enumerate}
\def\labelenumi{\alph{enumi})}
\setcounter{enumi}{1}
\tightlist
\item
  There exists a unique positive sesquilinear form \(b_{\mu}\) on C{[}X{]} such that \(b_{\mu}(X^{i}, X^{j}) = \mu_{i+j}\) for all integers i and j.
\end{enumerate}

We denote \(E_{\mu}\) the separated completion Hilbert space of C{[}X{]} for the form \(b_{\mu}\) . c) The linear map \(u_{\mu}\) from C{[}X{]} to C{[}X{]} defined by \(u_{\mu}(P) = XP\) induces, by passing to quotients, a symmetric partial operator \(u_{\mu}\) on the Hilbert space \(E_{\mu}\) .

\begin{enumerate}
\def\labelenumi{\alph{enumi})}
\setcounter{enumi}{3}
\item
  The partial operator \(u_{\mu}\) admits a self-adjoint extension.
\item
  There exists a measure \(\nu\in\mathcal{M}_{\mathrm{m}}(\mathbf{R})\) such that \(\mu_{n}=\nu(x^{n})\) for all \(n\in N\) . f) The sesquilinear form \(b_{\mu}\) is separating if and only if, for every measure \(\nu\in\mathcal{M}_{\mathrm{m}}(\mathbf{R})\) such that \(\mu_{n}=\nu(x^{n})\) for all \(n\in N\) , the support of \(\nu\) is infinite. (Observe that the left-hand side of inequality (24) is equal to \(b_{\mu}(p,p)\) for \(p=a_{0}+\cdots+a_{n}X^{n}\) .)
\end{enumerate}

\begin{enumerate}
\def\labelenumi{\arabic{enumi})}
\setcounter{enumi}{21}
\tightlist
\item
  We resume the notation of the preceding exercise.
\end{enumerate}

Let \(\boldsymbol{\mu} = (\mu_n)_{n\in \mathbf{N}}\) a sequence of real numbers satisfying conditions (24). We consider the question of uniqueness of a measure \(\nu \in \mathcal{M}_{\mathrm{m}}(\mathbf{R})\) such that \(\mu_{n} = \nu (x^{n})\) for all \(n\in \mathbf{N}\) . We assume that the sesquilinear form \(b_{\boldsymbol{\mu}}\) is separating (cf. question \(f\) ) of the preceding exercise).

\begin{enumerate}
\def\labelenumi{\alph{enumi})}
\item
  If there exists a measure \(\nu \in \mathcal{M}_{\mathrm{m}}(\mathbf{R})\) with compact support such that \(\nu(x^{n}) = \mu_{n}\) for all \(n \in \mathbf{N}\) , then \(\nu\) is the unique measure in \(\mathcal{M}_{\mathrm{m}}(\mathbf{R})\) such that \(\nu(x^{n}) = \mu_{n}\) for all \(n \in \mathbf{N}\) .
\item
  Let \(\nu \in \mathcal{M}_{\mathrm{m}}(\mathbf{R})\) such that \(\nu (x^n) = \mu_n\) for all \(n\in \mathbf{N}\) . For every \(p\in \mathbf{C}[\mathrm{X}]\) , we have \(p\in \mathcal{L}^2 (\mathbf{R},\nu)\) and \(\| p\| = b_{\mu}(p,p)\) .
\item
  The canonical injection of C{[}X{]} in \(\mathcal{L}^{2}(\mathbf{R},\nu)\) defines by continuity and by passing to the quotient an isometry of \(E_{\mu}\) in \(\mathrm{L}^{2}(\mathbf{R},\nu)\) In what follows, we identify \(E_{\mu}\) with its image, which is a closed subspace of \(\mathrm{L}^{2}(\mathbf{R},\nu)\) .
\end{enumerate}

\(d)\) One has \(u_{\mu} \subset m_x\) , where \(m_x\) denotes the operator of multiplication by \(x\) in \(\mathrm{L}^2(\mathbf{R}, \nu)\) .

\begin{enumerate}
\def\labelenumi{\alph{enumi})}
\setcounter{enumi}{4}
\item
  The measure \(\nu\) is the spectral measure of the constant function 1 relative to \(m_x\) .
\item
  Let \(\nu \in \mathcal{M}_{\mathrm{m}}(\mathbf{R})\) . We have \(\nu(x^n) = \mu_n\) for all \(n \in \mathbf{N}\) if and only if there exists a Hilbert space F and a partial self-adjoint operator \(v\) on F satisfying the following conditions:
\end{enumerate}

\begin{enumerate}
\def\labelenumi{(\roman{enumi})}
\item
  The space \(E_{\mu}\) is a closed subspace of F, and the canonical injection from \(E_{\mu}\) in F is an isometry;
\item
  One has \(u_{\mu} \subset v\) ;
\item
  The measure \(\nu\) is the spectral measure of the constant function \(1 \in E_{\mu}\) relative to v.
\end{enumerate}

\begin{enumerate}
\def\labelenumi{\alph{enumi})}
\setcounter{enumi}{6}
\item
  There exists a unique measure \(\nu \in \mathcal{M}_{\mathrm{m}}(\mathbf{R})\) such that \(\nu(x^{m}) = \mu_{n}\) for all \(n \in \mathbf{N}\) if and only if the partial operator \(u_{\boldsymbol{\mu}}\) is essentially self-adjoint.
\item
  Let \(\nu_{lg}\) the positive measure on R with support in \(R_{+}\) whose restriction to \(R_{+}^{*}\) to the density
\end{enumerate}

\[
t \mapsto \frac {1}{t} \exp (- \log (t) ^ {2} / 2)
\]

with respect to Lebesgue measure. Let \(\widetilde{\nu}_{lg}\) the positive measure on R with support in \(R_{+}\) whose restriction to \(R_{+}^{*}\) to the density \(t \mapsto 1 + \sin(2\pi \log(t))\) with respect to \(\nu_{lg}\) .

One has \(\nu_{\mathrm{lg}} \in \mathcal{M}_{\mathrm{m}}(\mathbf{R})\) and \(\widetilde{\nu}_{\mathrm{lg}} \in \mathcal{M}_{\mathrm{m}}(\mathbf{R})\) .

\begin{enumerate}
\def\labelenumi{\roman{enumi})}
\tightlist
\item
  One has \(\nu_{\mathrm{lg}}(x^n) = \widetilde{\nu}_{\mathrm{lg}}(x^n)\) for all \(n\in \mathbf{N}\) .
\end{enumerate}

\begin{enumerate}
\def\labelenumi{\arabic{enumi})}
\setcounter{enumi}{22}
\tightlist
\item
  We resume the notations of the previous exercise; in particular, we consider a sequence \(\boldsymbol{\mu} = (\mu_n)_{n\in \mathbf{N}}\) of real numbers satisfying conditions (24) and such that the sesquilinear form \(b_{\mu}\) is separating. We identify \(\mathbf{C}[\mathrm{X}]\) with a dense subspace of \(\mathrm{E}_{\mu}\) .
\end{enumerate}

\begin{enumerate}
\def\labelenumi{\alph{enumi})}
\tightlist
\item
  There exists a unique sequence \((p_{n})_{n\in\mathbf{N}}\) of polynomials in C{[}X{]} such that the following conditions are satisfied (“orthogonal polynomials of the first kind associated with \(\mu\) »):
\end{enumerate}

\begin{enumerate}
\def\labelenumi{(\roman{enumi})}
\item
  The family \((p_n)\) is an orthonormal basis of \(\mathrm{E}_{\mu}\) ;
\item
  The polynomial \(p_n\) is of degree \(n\) for all \(n \in \mathbf{N}\) ;
\item
  The coefficient of \(\mathrm{X}^n\) of \(p_n\) belongs to \(\mathbf{R}_+^*\) .
\end{enumerate}

Moreover, we have \(p_n \in \mathbf{R}[\mathbf{X}]\) for all \(n \in \mathbf{N}\) and \(p_0 = 1 / \sqrt{\mu_0}\) . We shall also denote \(p_{-1} = 0\) .

\begin{enumerate}
\def\labelenumi{\alph{enumi})}
\setcounter{enumi}{1}
\tightlist
\item
  There exist sequences \((a_{n})_{n\geqslant -1}\) in \(\mathbf{R}_{+}^{*}\) and \((b_{n})_{n\in \mathbf{N}}\) in \(\mathbf{R}\) such that \(a_{-1} = 1\) and
\end{enumerate}

\[
u _ {\pmb {\mu}} (p _ {n}) = a _ {n} p _ {n + 1} + b _ {n} p _ {n} + a _ {n - 1} p _ {n - 1}
\]

for all \(n \in \mathbf{N}\) . These sequences are unique.

\begin{enumerate}
\def\labelenumi{\alph{enumi})}
\setcounter{enumi}{2}
\item
  For every \(n \geqslant 1\) , the coefficient of \(X^{n}\) of the orthogonal polynomial \(p_{n}\) is equal to \((a_{0} \cdots a_{n-1})^{-1}\) .
\item
  Let \(\mathcal{F}\) the complex vector space of sequences \((z_{n})_{n\geqslant -1}\) of complex numbers such that \(z_{-1} = 0\) . We denote \(\widetilde{v}_{\mu}\) the linear map from \(\mathcal{F}\) into itself such that \(\widetilde{v}_{\mu}(z_n) = (w_n)_{n\in \mathbf{N}}\) where \(w_{-1} = 0\) and
\end{enumerate}

\[
w _ {n} = a _ {n} z _ {n + 1} + b _ {n} z _ {n} + a _ {n - 1} z _ {n - 1}
\]

for all \(n \in N\) .

For every \(s \in \mathbf{C}\) the sequence \(z_s = (p_n(s))_{n \geqslant -1}\) is the unique sequence in \(\mathcal{F}\) such that \((z_s)_0 = p_0(z)\) and \(\widetilde{v}_{\boldsymbol{\mu}}(z_s) = sz_s\) .

\begin{enumerate}
\def\labelenumi{\alph{enumi})}
\setcounter{enumi}{4}
\tightlist
\item
  Let \(z\) and \(w\) be elements of \(\mathcal{F}\) . One defines the sequence \(\mathrm{W}(z,w) = (\mathrm{W}_n(z,w))_{n\in \mathbf{N}}\) by setting
\end{enumerate}

\[
\mathrm{W} _ {n} (z, w) = a _ {n} (z _ {n + 1} w _ {n} - z _ {n} w _ {n + 1})
\]

for \(n\in \mathbf{N}\) . We have

\[
\mathrm{W} _ {n} (z, w) = \sum_ {j = 0} ^ {n} (\widetilde {v} _ {\boldsymbol {\mu}} (z) _ {j} w _ {j} - z _ {j} \widetilde {v} _ {\boldsymbol {\mu}} (w) _ {j})
\]

for all \(n \in N\) .

\(f)\) Let \(\mathrm{P} \subset \ell^2(\mathbf{N})\) the space of sequences with finite support. For every \(z \in \mathrm{P}\) the sequence \(\widetilde{v}_{\boldsymbol{\mu}}(z)\) belongs to \(\ell^2(\mathbf{N})\) .

We denote \(v_{\mu}\) the partial operator on \(\ell^2(\mathbf{N})\) of domain P given by \(z \mapsto \widetilde{v}_{\mu}(z)\) . It is symmetric.

\begin{enumerate}
\def\labelenumi{\alph{enumi})}
\setcounter{enumi}{6}
\tightlist
\item
  The adjoint of \(v_{\mu}\) is the partial operator with domain
\end{enumerate}

\[
\mathrm{dom} (v _ {\boldsymbol {\mu}} ^ {*}) = \{z \in \ell^ {2} (\mathbf {N}) | \widetilde {v} _ {\boldsymbol {\mu}} (z) \in \ell^ {2} (\mathbf {N}) \}
\]

such that \(v_{\mu}^{*}(z) = \widetilde{v}_{\mu}(z)\) for \(z \in \mathrm{dom}(v_{\mu}^{*})\) .

\(h)\) For all \(z\) and \(w\) in \(\mathrm{dom}(v_{\mu}^{*})\) , the limit of \(\mathrm{W}_{n}(\overline{z}, w)\) as \(n \to +\infty\) exists and is equal to \(\langle w | v_{\mu}^{*}(z) \rangle - \langle v_{\mu}^{*}(w) | z \rangle\) .

\begin{enumerate}
\def\labelenumi{\roman{enumi})}
\tightlist
\item
  Let \(s \in C\) . The kernel of \(v_{\boldsymbol{\mu}}^{*} - s \cdot 1_{\ell^{2}(\mathbf{N})}\) has dimension \(\leqslant 1\) ; it has dimension 1 if and only if \(z_{s} \in \ell^{2}(\mathbf{N})\) , and it is then generated by \(z_{s}\) .
\end{enumerate}

\begin{enumerate}
\def\labelenumi{\alph{enumi})}
\setcounter{enumi}{9}
\tightlist
\item
  The closure of the partial operator \(v_{\mu}\) (resp. of \(u_{\mu}\) ) has deficiency index (0, 0) or (1, 1); the first case holds if and only if there exists \(s \in \mathbf{C} - \mathbf{R}\) such that \(z_s \notin \ell^2(\mathbf{N})\) . (For the case of \(u_{\mu}\) , observe that the map \(X^n \mapsto e_n\) , where \((e_n)\) is the canonical basis of \(\ell^2(\mathbf{N})\) , induces an isometric isomorphism of \(E_{\mu}\) in \(\ell^2(\mathbf{N})\) which makes it possible to identify \(u_{\mu}\) with \(v_{\mu}\) .)
\end{enumerate}

\(k)\) Let \(s \in \mathbf{C}\) and \(f \in \mathrm{Ker}(u_{\mu}^{*} - s \cdot 1_{\mathrm{E}_{\mu}})\) . If \(\langle f | 1 \rangle = 0\) , then \(f = 0\) .

\begin{enumerate}
\def\labelenumi{\alph{enumi})}
\setcounter{enumi}{11}
\item
  Let \(u_1\) and \(u_2\) of the self-adjoint extensions of \(u_\mu\) onto \(\mathrm{E}_\mu\) . We have \(u_1 = u_2\) if and only if there exists \(s \in \mathbf{C} - \mathbf{R}\) such that \(\langle \mathrm{R}(u_1, s)1|1\rangle = \langle \mathrm{R}(u_2, s)1|1\rangle\) . (Suppose that such an equality holds; set \(f_1 = \mathrm{R}(u_1, s)1\) , and show that \(f_1\) does not belong to the domain of \(\overline{u}_\mu\) ; deduce that \(u_1\) coincides with the closure of the restriction of \(u_1\) to \(\mathrm{dom}(\overline{u}_\mu) \oplus \mathbf{C}f_1\) . Reason similarly with \(u_2\) and \(f_2 = \mathrm{R}(u_2, s)1\) , then prove that \(f_1 = f_2\) to conclude that \(u_1 = u_2\) .) m) Suppose that \(u_\mu\) is not essentially self-adjoint. There then exist at least two measures \(\nu \in \mathcal{M}_{\mathrm{m}}(\mathbf{R})\) satisfying \(\nu(x^n) = \mu_n\) for all \(n \in \mathbf{N}\) . n) In what follows, assume that \(u_\mu\) is essentially self-adjoint. For every \(s \in \mathbf{C} - \mathbf{R}\) , there exists a sequence \((r_{s,m})_{m \in \mathbf{N}}\) of polynomials such that \((x - s)r_{s,m}\) converges in \(\mathrm{E}_\mu\) to the constant function 1 as \(m \to +\infty\) .
\item
  Let \(\nu\in\mathcal{M}_{\mathrm{m}}(\mathbf{R})\) a measure such that \(\nu(x^{n})=\mu_{n}\) for all \(n\in\mathbf{N}\) . For every \(m\in\mathbf{N}\) , the inequality
\end{enumerate}

\[
\left| \int_ {\mathbf {R}} \frac {d \nu (t)}{t - s} - \int_ {\mathbf {R}} r _ {s, m} (t) d \nu (t) \right| \leqslant \frac {\mu_ {0} ^ {2}}{| \mathcal {I} (s) ^ {2} |} \| 1 - (x - s) r _ {s, m} \| _ {\mathrm{E} _ {\mu}}
\]

is valid.

\begin{enumerate}
\def\labelenumi{\alph{enumi})}
\setcounter{enumi}{15}
\tightlist
\item
  The measure \(\nu\) is unique.
\end{enumerate}

\begin{enumerate}
\def\labelenumi{\arabic{enumi})}
\setcounter{enumi}{23}
\tightlist
\item
  We resume the notations of the preceding exercises; in particular, we consider a sequence \(\boldsymbol{\mu} = (\mu_n)_{n\in \mathbf{N}}\) of real numbers satisfying conditions (24) and such that the sesquilinear form \(b_{\boldsymbol{\mu}}\) is separating.
\end{enumerate}

\begin{enumerate}
\def\labelenumi{\alph{enumi})}
\tightlist
\item
  There exists a unique sequence of polynomials \((q_{n})_{n\in\mathbf{N}}\) with real coefficients such that for every \(s\in C\) , the formula
\end{enumerate}

\[
q _ {n} (s) = b _ {\pmb {\mu}} \Big (1, \frac {p _ {n} - p _ {n} (s)}{\mathrm{X} - s} \Big)
\]

is valid (“orthogonal polynomials of the second kind”).

\begin{enumerate}
\def\labelenumi{\alph{enumi})}
\setcounter{enumi}{1}
\tightlist
\item
  Let \(\nu \in \mathcal{M}_{\mathrm{m}}(\mathbf{R})\) a measure on \(\mathbf{R}\) such that \(\nu(x^n) = \mu_n\) for all \(n \in \mathbf{N}\) . For every \(s \in \mathbf{C} - \mathbf{R}\) , we have
\end{enumerate}

\[
\int_ {\mathbf {R}} \frac {p _ {n} (t)}{t - s} d \nu (t) = q _ {n} (s) + p _ {n} (s) \int_ {\mathbf {R}} \frac {d \nu (t)}{t - s}.
\]

\begin{enumerate}
\def\labelenumi{\alph{enumi})}
\setcounter{enumi}{2}
\item
  Let \(s \in C-R\) . The sequence \((p_{n}(s))_{n \in \mathbf{N}}\) belongs to \(\ell^{2}(\mathbf{N})\) if and only if the sequence \((q_{n}(s))_{n \in \mathbf{N}}\) belongs to \(\ell^{2}(\mathbf{N})\) . (Let \(\sigma(s) = \int(t-s)^{-1} d\nu(t)\) ; using the preceding question, prove that the sequence \((q_{n}(s) + \sigma(s)p_{n}(s))_{n \in \mathbf{N}}\) belongs to \(\ell^{2}(\mathbf{N})\) .)
\item
  For every \(n \in \mathbf{N}\) , the formula
\end{enumerate}

\[
a _ {n} (p _ {n + 1} (\mathrm{X}) p _ {n} (\mathrm{Y}) - p _ {n} (\mathrm{X}) p _ {n + 1} (\mathrm{Y})) = (\mathrm{X} - \mathrm{Y}) \sum_ {j = 0} ^ {n} p _ {j} (\mathrm{X}) p _ {j} (\mathrm{Y})
\]

is valid in \(\mathbf{C}[\mathrm{X},\mathrm{Y}]\) .

\begin{enumerate}
\def\labelenumi{\alph{enumi})}
\setcounter{enumi}{4}
\tightlist
\item
  For every \(n \in \mathbf{N}\) , we have the relation
\end{enumerate}

\[
p _ {n} q _ {n + 1} - p _ {n + 1} q _ {n} = \frac {1}{a _ {n}}.
\]

\begin{enumerate}
\def\labelenumi{\alph{enumi})}
\setcounter{enumi}{5}
\tightlist
\item
  If the series with general term \(a_{n}^{-1}\) is divergent, then the measure \(\nu\) is the unique measure such that \(\nu(x^{n}) = \mu_{n}\) for all \(n \in N\) .
\end{enumerate}

\begin{enumerate}
\def\labelenumi{\arabic{enumi})}
\setcounter{enumi}{24}
\tightlist
\item
  We resume the notations of the preceding exercises; in particular, we consider a sequence \(\boldsymbol{\mu} = (\mu_{n})_{n \in \mathbf{N}}\) of real numbers satisfying conditions (24) and such that the sesquilinear form \(b_{\mu}\) is separating.
\end{enumerate}

Suppose that \(\mu_{0}=1\) and such that the Carleman condition is satisfied, that is, the series with general term \(\mu_{2n}^{-1/(2n)}\) for \(n\geqslant1\) is divergent. a) We have

\[
\mu_ {2 n} ^ {- 1 / (2 n)} \leqslant (a _ {0} \dots a _ {n - 1}) ^ {- 1 / n} \leqslant \frac {e}{n ^ {2}} \sum_ {j = 1} ^ {n} \frac {j}{a _ {j - 1}}
\]

for every integer \(n \geqslant 1\) . (To establish the second inequality, one may prove and use the inequality \((n/e)^{n} \leqslant n!\) for \(n \geqslant 1\) .)

\begin{enumerate}
\def\labelenumi{\alph{enumi})}
\setcounter{enumi}{1}
\item
  The series with general term \(a_{n}^{-1}\) for \(n \geqslant 1\) is divergent.
\item
  There exists a unique measure \(\nu \in \mathcal{M}_{\mathrm{m}}(\mathbf{R})\) such that \(\nu(x^n) = \mu_n\) for all \(n \in \mathbf{N}\) .
\end{enumerate}

\begin{enumerate}
\def\labelenumi{\arabic{enumi})}
\setcounter{enumi}{25}
\tightlist
\item
  We denote \(x\) the identity function of \(\mathbf{R}_{+}\) .
\end{enumerate}

\begin{enumerate}
\def\labelenumi{\alph{enumi})}
\tightlist
\item
  Let \(\nu\) a positive measure on \(R_{+}\) such that \(x^{n}\) is \(\nu\) -integrable for every integer \(n \geqslant 0\) . Denote by \(\mu_{n} = \nu(x^{n})\) . For every integer \(n \geqslant 0\) and for all complex numbers \(a_{0}, \ldots, a_{n}\) , we have
\end{enumerate}

\[
\sum_ {i = 0} ^ {n} \sum_ {j = 0} ^ {n} \overline {{a}} _ {i} a _ {j} \mu_ {i + j} \geqslant 0
\]

\[
\sum_ {i = 0} ^ {n} \sum_ {j = 0} ^ {n} \overline {{a}} _ {i} a _ {j} \mu_ {i + j + 1} \geqslant 0.
\]

\begin{enumerate}
\def\labelenumi{\alph{enumi})}
\setcounter{enumi}{1}
\tightlist
\item
  If \((\mu_{n})_{n\geqslant0}\) is a sequence of real numbers satisfying the condition of the preceding question, then there exists a positive measure \(\nu\) onto \(R_{+}\) such that \(\mu_{n}=\nu(x^{n})\) for all \(n\in N\) . (“Stieltjes moment problem”; use the method of Exercise 21, appealing to the Friedrichs extension.)
\end{enumerate}

\begin{enumerate}
\def\labelenumi{\arabic{enumi})}
\setcounter{enumi}{26}
\tightlist
\item
  We denote \(M = X^{2}Y^{2}(X^{2} + Y^{2} - 3) \in \mathbf{R}[X, Y]\) (“Motzkin polynomial”).
\end{enumerate}

\begin{enumerate}
\def\labelenumi{\alph{enumi})}
\item
  We have M(x,y) ≥ 0 for all x and y in R.
\item
  There exists no finite family \((p_{i})_{i\in I}\) in R{[}X,Y{]} such that \(M=\sum p_{i}^{2}\) . (If such a family existed, show that each \(p_{i}\) must be of the form \(a+bXY\) with \(a\in R\) and \(b\in R[X,Y]\) of degree at most 1; then compute the coefficient of \((\mathrm{XY})^{2}\) .)
\item
  Let \(\tau\colon\mathbf{N}^{2}\to\mathbf{N}-\{0\}\) the map defined by
\end{enumerate}

\[
\tau (0, 0) = 1, \quad \tau (1, 2) = 2, \quad \tau (2, 1) = 3, \quad \tau (1, 1) = 4, \quad \tau (1, 0) = 5,
\]

\[
\tau (0, 1) = 6, \quad \tau (2, 0) = 7, \quad \tau (0, 2) = 8, \quad \tau (3, 0) = 9, \quad \tau (0, 3) = 1 0
\]

and

\[
\tau (i, j) = j + 1 + (i + j) (i + j + 1) / 2 \text {   if   } i + j \geqslant 4.
\]

The map \(\tau\) is a bijection.

\begin{enumerate}
\def\labelenumi{\alph{enumi})}
\setcounter{enumi}{3}
\tightlist
\item
  Let \((\nu_{n})_{n\geqslant 1}\) the sequence defined by
\end{enumerate}

\[
\nu_ {i} = 1 \text {if} i \in \{1, 2, 3 \}, \quad \nu_ {4} = 4, \quad \nu_ {n} = (n!) ^ {(n + 1)!} \text {if} n \geqslant 5.
\]

For \(n\) and \(m\) in \(\mathbf{N}\) we set \(\mu_{n,m} = 0\) if \(n\) or \(m\) is odd and \(\mu_{n,m} = \nu_{\tau(n/2,m/2)}\) otherwise. The linear map \(\alpha\) of \(\mathbf{C}[\mathrm{X},\mathrm{Y}]\) in \(\mathbf{C}\) such that \(\alpha(\mathrm{X}^n\mathrm{Y}^m) = \mu_{n,m}\) satisfies \(\alpha(p^2) \geqslant 0\) for all \(p \in \mathbf{R}[\mathrm{X},\mathrm{Y}]\) . (It suffices to verify that, for every integer \(k \in \mathbf{N}\) the determinant of the matrix \((a_{m,m})_{1 \leqslant n,m \leqslant k}\) is positive, where \(a_{n,m}\) is defined by \(a_{\tau(i,j),\tau(i',j')} = \alpha(\mathrm{X}^{i+i'}\mathrm{Y}^{j+j'})\) for all \((i,j,i',j')\) ; then proceed by induction on \(k\) .)

\begin{enumerate}
\def\labelenumi{\alph{enumi})}
\setcounter{enumi}{4}
\tightlist
\item
  There exists no positive measure \(\nu\) onto \(R^{2}\) such that \((x,y)\mapsto x^{n}y^{m}\) let \(\mu\) that is integrable for all \((n,m)\in\mathbf{N}^{2}\) , and such that
\end{enumerate}

\[
\int_ {\mathbf {R} ^ {2}} x ^ {n} y ^ {m} d \nu (x, y) = \mu_ {n, m}
\]

for all \((n,m)\in\mathbf{N}^{2}\) .

\begin{enumerate}
\def\labelenumi{\alph{enumi})}
\setcounter{enumi}{5}
\tightlist
\item
  There exists a Hilbert space E and a densely defined partial operator u on E satisfying the following conditions:
\end{enumerate}

\begin{enumerate}
\def\labelenumi{(\roman{enumi})}
\item
  One has \(\operatorname{dom}(u) \subset \operatorname{dom}(u^*)\) ;
\item
  One has \(\|u(x)\|=\|u^{*}(x)\|\) for all \(x\in\mathrm{dom}(u)\) ;
\item
  There does not exist a Hilbert space F containing E as a subspace, such that the canonical injection of E into F is isometric, and a partial normal operator v on F such that \(u \subset v\) .
\end{enumerate}

Compare this result with exercise 24 of IV, p. 352.

\begin{enumerate}
\def\labelenumi{\arabic{enumi})}
\setcounter{enumi}{27}
\tightlist
\item
  For every integer \(k \geqslant 0\) , we denote \(b_k\) the number of equivalence relations on the set \(\{1, \ldots, k\}\) .
\end{enumerate}

For every \(\lambda > 0\) , we denote \(\varpi_{\lambda}\) the Poisson measure with parameter \(\lambda\) onto \(\mathbf{R}\) (cf. exercise 41 of II, p. 282), that is, the measure of total mass 1 supported on \(\mathbf{N}\) such that \(\varpi_{\lambda}(j) = e^{-\lambda} \lambda^{j} / j!\) for all \(j \in \mathbf{N}\) .

\begin{enumerate}
\def\labelenumi{\alph{enumi})}
\tightlist
\item
  The relation
\end{enumerate}

\[
b _ {k + 1} = \sum_ {j = 0} ^ {k} {\binom {k} {j}} b _ {j}
\]

is valid for every integer \(k \in N\) .

\begin{enumerate}
\def\labelenumi{\alph{enumi})}
\setcounter{enumi}{1}
\tightlist
\item
  For any function \(f: N \to C\) such that the function \(x \mapsto xf(x)\) is \(\varpi\) -integrable, we have
\end{enumerate}

\[
\int_ {\mathbf {R}} x f (x) d \varpi_ {\lambda} (x) = \lambda \int_ {\mathbf {R}} f (x + 1) d \varpi_ {\lambda} (x)
\]

(“Chen’s formula”).

Additionally, if \(\mu\) is a bounded measure on R satisfying this property, then there exists \(c \in R\) such that \(\mu = c\varpi_{\lambda}\) .

\begin{enumerate}
\def\labelenumi{\alph{enumi})}
\setcounter{enumi}{2}
\tightlist
\item
  For every \(k \in \mathbf{N}\) , the function \(x \mapsto x^k\) is \(\varpi_1\) -integrable and we have
\end{enumerate}

\[
\varpi_ {1} (x \mapsto x ^ {k}) = b _ {k}.
\]

Additionally, the measure \(\varpi_{1}\) is the unique positive measure on R satisfying these equations.

\begin{enumerate}
\def\labelenumi{\alph{enumi})}
\setcounter{enumi}{3}
\tightlist
\item
  For every integer \(k \geqslant 0\) , we have
\end{enumerate}

\[
b _ {k} = \frac {1}{e} \sum_ {j \geqslant 0} \frac {j ^ {n}}{j !}
\]

(« Dobiński's formula ») and

\[
b _ {k} = \frac {2 k !}{e \pi} \mathcal {I} \left(\int_ {0} ^ {\pi} e ^ {e ^ {e ^ {i t}}} \sin (k t) d t\right)
\]

(“Cesàro formula”).

\begin{enumerate}
\def\labelenumi{\alph{enumi})}
\setcounter{enumi}{4}
\tightlist
\item
  For every integer \(n \geqslant 1\) , we denote \(\mu_{n}\) the positive measure on R that is the image of the normalized Haar measure of the symmetric group \(S_{n}\) by the mapping \(f_{n}\) of \(S_{n}\) in R which to \(\sigma \in S_{n}\) associates the number of fixed points of \(\sigma\) Let \(k \in N\) . For every integer \(n \geqslant k\) , we have
\end{enumerate}

\[
\mu_ {n} (x \mapsto x ^ {k}) = b _ {k}.
\]

\begin{enumerate}
\def\labelenumi{\alph{enumi})}
\setcounter{enumi}{5}
\tightlist
\item
  The sequence \((\mu_{n})\) converges weakly when \(n \to +\infty\) toward the Poisson measure \(\varpi_{1}\) .
\end{enumerate}

\begin{enumerate}
\def\labelenumi{\arabic{enumi})}
\setcounter{enumi}{28}
\tightlist
\item
  Let E be a complex Hilbert space. We denote by \(\mathcal{T}_b\) (resp. \(\mathcal{T}_s\) , \(\mathcal{T}_f\) ) the Banach space topology of \(\mathcal{L}(\mathrm{E})\) (resp. the topology of pointwise convergence, the locally convex topology defined by the seminorms \(u \mapsto \langle x \mid u(y) \rangle\) for \((x, y) \in \mathrm{E} \times \mathrm{E}\) ).
\end{enumerate}

\begin{enumerate}
\def\labelenumi{\alph{enumi})}
\item
  Let \(\lambda_0\in \mathbf{C} - \mathbf{R}\) The topology \(\mathcal{T}_b\) -resolvent (resp. \(\mathcal{T}_s\) -resolvent) is the coarsest topology on \(\mathcal{A}(\mathrm{E})\) such that the map \(u\mapsto \mathrm{R}(u,\lambda_0)\) of \(\mathcal{A}(\mathrm{E})\) in \(\mathcal{L}(\mathrm{E})\) endowed with the topology \(\mathcal{T}_b\) (resp. \(\mathcal{T}_s\) ) is continuous.
\item
  The topology \(T_{f}\) -resolvent coincides with the topology \(T_{s}\) -resolvent.
\item
  Let \(X \subset \mathcal{L}(E)\) a bounded set of Hermitian endomorphisms of E. The topology induced on X by the topology \(T_{b}\) -resolvent (resp. the topology \(T_{s}\) -resolvent) coincides with the topology induced by the topology \(T_{b}\) (esp. the topology \(T_{s}\) ).
\item
  The property from part (a) does not hold for the topology. \(T_{f}\) -resolvent.
\item
  Give an example of a sequence. \((u_{n})\) of partial self-adjoint operators that converges for the topology \(T_{b}\) -resolvent and function \(f \in \mathcal{C}_{b}(\mathbf{R})\) such that \(f(u_{n})\) does not converge to \(f(u)\) in \(\mathcal{L}(\mathrm{E})\) .
\end{enumerate}

\(f\) Let \((u_n)\) a sequence of partial self-adjoint operators on E and \(u\) a self-adjoint partial operator on E. If there exists a dense subspace F of \(\mathrm{E}_u\) contained in \(\mathrm{dom}(u_n)\) for all \(n\) such that \(u_n(x)\) converges to \(u(x)\) for all \(x \in \mathrm{F}\) , then \((u_n)\) converges to \(u\) for the topology \(\mathcal{T}_s\) -resolvent.

\begin{enumerate}
\def\labelenumi{\arabic{enumi})}
\setcounter{enumi}{29}
\tightlist
\item
  Let E be a complex Hilbert space. Denote by \(T_{b}\) the topology of bounded convergence on \(\mathcal{L}(\mathrm{E})\) .
\end{enumerate}

\begin{enumerate}
\def\labelenumi{\alph{enumi})}
\tightlist
\item
  Let \(u\) a self-adjoint partial operator on E. For every \(\lambda \in \mathbf{C} - \mathbf{R}\) such that the imaginary part of \(\lambda\) is negative, we have
\end{enumerate}

\[
\mathrm{R} (u, \lambda) = i \int_ {0} ^ {\infty} e ^ {- i t \lambda} e ^ {i t u} d t.
\]

\begin{enumerate}
\def\labelenumi{\alph{enumi})}
\setcounter{enumi}{1}
\tightlist
\item
  The topology \(T_{b}\) -resolvent is the coarsest topology on \(\mathcal{A}(E)\) such that the maps \(u \mapsto e^{itu}\) of \(\mathcal{A}(E)\) in \(\mathcal{L}(E)\) are continuous for every \(t \in R\) .
\end{enumerate}

\begin{enumerate}
\def\labelenumi{\arabic{enumi})}
\setcounter{enumi}{30}
\tightlist
\item
  Let E be a complex Hilbert space. Denote by \(T_{s}\) the topology of pointwise convergence \(\mathcal{L}(\mathrm{E})\) Let H denote the set of complex numbers with strictly positive imaginary part.
\end{enumerate}

Let \((u_{n})\) a sequence in \(\mathcal{A}(\mathrm{E})\) . One assumes that

\begin{enumerate}
\def\labelenumi{(\roman{enumi})}
\item
  There exists \(\lambda_{+} \in \mathbf{H}\) such that \(\mathrm{R}(u_n, \lambda_{+})\) converges in \(\mathcal{L}(\mathrm{E})\) endowed with the topology \(\mathcal{T}_s\) ; we denote by \(v_{+}\) its limit;
\item
  There exists \(\lambda_{-} \in \overline{\mathbf{H}}\) such that \(\mathrm{R}(u_n, \lambda_{-})\) converges in \(\mathcal{L}(\mathrm{E})\) endowed with the topology \(\mathcal{T}_s\) ;
\item
  The image F of the endomorphism \(v_{+}\) is dense in E.
\end{enumerate}

Let D be the open disk in C consisting of \(\lambda \in H\) such that

\[
| \lambda - \lambda_ {+} | <   \frac {1}{| \mathcal {I} (\lambda_ {+}) |}.
\]

\begin{enumerate}
\def\labelenumi{\alph{enumi})}
\tightlist
\item
  For every \(\lambda \in D\) , the series
\end{enumerate}

\[
r _ {\lambda} = \sum_ {n = 0} ^ {\infty} (\lambda_ {+} - \lambda) v _ {+} ^ {n + 1}
\]

converges in \(\mathcal{L}(\mathrm{E})\) . Moreover, for every \(\lambda \in \mathrm{D}\) the sequence \((\mathrm{R}(u_n,\lambda))_n\) converges to \(r_{\lambda}\) in \(\mathcal{L}(\mathrm{E})\) endowed with the topology \(\mathcal{T}_s\) .

\begin{enumerate}
\def\labelenumi{\alph{enumi})}
\setcounter{enumi}{1}
\item
  There exists a unique holomorphic map from H to \(\mathcal{L}(\mathrm{E})\) which coincides on D with the map \(\lambda\mapsto r_{\lambda}\) We still denote it by \(\lambda\mapsto r_{\lambda}\) .
\item
  There exists a unique holomorphic map from \(\mathbf{C} - \mathbf{R}\) in \(\mathcal{L}(\mathrm{E})\) which coincides on \(\mathbf{H}\) with the map \(\lambda \mapsto r_{\lambda}\) and which satisfies \(r_{\lambda}^{*} = r_{\overline{\lambda}}\) for all \(\lambda \in \mathbf{C} - \mathbf{R}\) It satisfies \(r_{\lambda_{-}} = v_{-}\) , and for all \(\lambda \in \mathbf{C} - \mathbf{R}\) the sequence \((\mathrm{R}(u_n,\lambda))\) converges to \(r_{\lambda}\) in \(\mathcal{L}(\mathrm{E})\) endowed with the topology \(\mathcal{T}_s\) .
\item
  For all \(\lambda\) and \(\mu\) in \(\mathbf{C} - \mathbf{R}\) , we have \(r_{\lambda} - r_{\mu} = (\mu - \lambda) = r_{\mu}r_{\lambda}\) and \(r_{\lambda}r_{\mu} = r_{\mu}r_{\lambda}\) .\\
\item
  The image of \(r_{\lambda}\) is equal to F for all \(\lambda \in \mathbf{C} - \mathbf{R}\) .
\end{enumerate}

\(f)\) For every \(\lambda\in\mathbf{C}\text{—}\mathbf{R}\) , the endomorphism \(r_{\lambda}\) is injective; the partial operator \(\lambda1_{\mathrm{E}}-r_{\lambda}^{-1}\) of domain F does not depend on the choice of \(\lambda\in\mathbf{C}\text{—}\mathbf{R}\) ; we denote it \(u\) . \(g)\) The partial operator \(u\) is self-adjoint; for all \(\lambda\in\mathbf{C}\text{—}\mathbf{R}\) , we have R( \(u,\lambda\) ) = \(r_{\lambda}\) .

\begin{enumerate}
\def\labelenumi{\alph{enumi})}
\setcounter{enumi}{7}
\tightlist
\item
  The sequence \((u_{n})\) converges to u for the topology \(T_{s}\) -resolvent. (“Kato–Trotter theorem”)
\end{enumerate}

\begin{enumerate}
\def\labelenumi{\arabic{enumi})}
\setcounter{enumi}{31}
\tightlist
\item
  Let E be a complex Hilbert space. Let u and v be partial self-adjoint operators on E such that \(u + v\) is self-adjoint.
\end{enumerate}

\begin{enumerate}
\def\labelenumi{\alph{enumi})}
\item
  For \(t \in \mathbf{R}^*\) , let \(v(t) = t^{-1}(e^{itu}e^{itv} - e^{it(u + v)}) \in \mathcal{L}(\mathrm{E})\) , and \(w(t)\) the restriction of \(v(t)\) to the Hilbert space \(\mathrm{E}_{u + v}\) . Let B be the set of continuous linear maps \(w(t)\) . The set B is bounded in the Banach space \(\mathcal{L}(\mathrm{E}_{u + v},\mathrm{E})\) . (Use the Banach-Steinhaus theorem.)
\item
  One has \(w(t) \to 0\) as \(t \to 0\) in \(R^{*}\) , into the space \(\mathcal{L}(\mathrm{E}_{u+v},\mathrm{E})\) endowed with the topology of compact convergence.
\item
  Let \(x\in \mathrm{E}_{u + v}\) . We have
\end{enumerate}

\[
\lim _ {n \to + \infty} (e ^ {i t u / n} e ^ {i t v / n}) ^ {n} x = e ^ {i t (u + v)} x.
\]

For \(n\in \mathbf{N}\) , write

\[
\begin{array}{l} (e ^ {i t u / n} e ^ {i t v / n}) ^ {n} - e ^ {i t (u + v)} = \\ \sum_ {i = 0} ^ {n - 1} (e ^ {i t u / n} e ^ {i t v / n}) ^ {k} \Big (e ^ {i t u / n} e ^ {i t v / n} - e ^ {i t (u + v) / n} \Big) (e ^ {i t (u + v) / n}) ^ {n - 1 - i} \end{array}
\]

and apply the result of the previous question to the compact image of the map \(t \mapsto e^{it(u+v)}x\) of \([-1,1]\) in \(E_{u+v}\) .

\begin{enumerate}
\def\labelenumi{\arabic{enumi})}
\setcounter{enumi}{32}
\tightlist
\item
  Let \(n \in \mathbf{N}\) . We denote \(\Delta\) the Laplacian on \(\mathrm{L}^2(\mathbf{R}^n)\) .
\end{enumerate}

\begin{enumerate}
\def\labelenumi{\alph{enumi})}
\tightlist
\item
  For every \(t \in R^{*}\) , the endomorphism \(e^{it\Delta}\) of \(\mathrm{L}^{2}(\mathbf{R}^{n})\) is given by
\end{enumerate}

\[
e ^ {i t \Delta} f (x) = \frac {1}{(2 i t) ^ {n / 2}} \int_ {\mathbf {R} ^ {n}} f (y) e ^ {i | x - y | ^ {2} / (4 t)} d y
\]

if \(f\in \mathcal{L}^1 (\mathbf{R}^n)\cap \mathcal{L}^2 (\mathbf{R}^n)\) (Use the Fourier transform.)

\begin{enumerate}
\def\labelenumi{\alph{enumi})}
\setcounter{enumi}{1}
\tightlist
\item
  Let C be the unitary operator on \(\mathrm{L}^{2}(\mathbf{R}^{n})\) obtained by passing to the quotient from the map \(f \mapsto \mathrm{C}(f)\) defined for \(f \in \mathcal{L}^{2}(\mathbf{R}^{n})\) by
\end{enumerate}

\[
\mathrm{C} (f) (x) = \frac {1}{(2 i t) ^ {n / 2}} e ^ {\frac {i x ^ {2}}{4 t}} \mathcal {F} f \Big (\frac {x}{2 t} \Big)
\]

for \(x\in \mathbf{R}^n\)

One has \(e^{it\Delta} \to C\) as \(t \to +\infty\) in the space \(\mathcal{L}(\mathrm{L}^{2}(\mathbf{R}^{n}))\) endowed with the topology of simple convergence.

\begin{enumerate}
\def\labelenumi{\alph{enumi})}
\setcounter{enumi}{2}
\tightlist
\item
  Suppose that \(1 \leqslant n \leqslant 3\) Let \(g \in \mathcal{L}^{\infty}(\mathbf{R}^{n}) + \mathcal{L}^{2}(\mathbf{R}^{n})\) with real values. Let u denote the partial operator \(u = \Delta + m_{g}\) ; it is self-adjoint on \(\mathrm{L}^{2}(\mathbf{R}^{n})\) (Exercise 14). For \(f \in \mathcal{L}^{1}(\mathbf{R}^{n}) \cap \mathcal{L}^{2}(\mathbf{R}^{n})\) and \(t \in R^{*}\) , we have
\end{enumerate}

\[
e ^ {i t u} f (x) = \lim _ {k \rightarrow + \infty} \frac {1}{(2 i t / k) ^ {k n / 2}} \int_ {\mathbf {R} ^ {n k}} f (y _ {k}) \exp \left(\frac {i t}{k} \sum_ {i = 1} ^ {k} \left(\frac {| y _ {i} - y _ {i - 1} | ^ {2}}{4 (t / k) ^ {2}} - g (y _ {i})\right) d y \right.
\]

where we write \(y = (y_{1}, \ldots, y_{k})\) with \(y_{i} \in R^{n}\) .

\begin{enumerate}
\def\labelenumi{\alph{enumi})}
\setcounter{enumi}{3}
\tightlist
\item
  Interpretation of the preceding formula.
\end{enumerate}

\begin{enumerate}
\def\labelenumi{\arabic{enumi})}
\setcounter{enumi}{33}
\tightlist
\item
  Let E be a complex Hilbert space. For all endomorphisms \(u\) and \(v\) of E, we denote \([u,v] = uv - vu\) and we define \(\mathrm{E}_{u,v} = \overline{\mathrm{Im}(u)} \cap \mathrm{Ker}(v)\) .
\end{enumerate}

Let \(p\) and \(q\) orthogonal projectors of E. We set \(a = p - q\) and \(b = 1_{\mathrm{E}} - p - q\) a) If there exists a unitary endomorphism \(u\) of E such that \(upu^{-1} = q\) and \(uqu^{-1} = p\) , then the Hilbert dimensions of \(\mathrm{E}_{p,q}\) and \(\mathrm{E}_{1 - p,1 - q}\) are equal.

\begin{enumerate}
\def\labelenumi{\alph{enumi})}
\setcounter{enumi}{1}
\tightlist
\item
  We have the relations
\end{enumerate}

\[
a ^ {2} + b ^ {2} = 1 _ {\mathrm{E}}, \qquad a b + b a = 0, \qquad [ p, a ] = [ q, a ] = [ p, b ] = [ q, b ] = 0.
\]

\begin{enumerate}
\def\labelenumi{\alph{enumi})}
\setcounter{enumi}{2}
\tightlist
\item
  Up to the question \(f\) ), included, we will assume that \(\mathrm{E}_{p,q} = \mathrm{E}_{1 - p,1 - q} = \{0\}\) . The endomorphism \(b\) is injective; let \(s\) the function on \(\mathbf{R}\) in \(\mathbf{R}\) such that
\end{enumerate}

\[
s (x) = \left\{ \begin{array}{l l} 1 & \text { if } x > 0 \\ 0 & \text { if } x = 0 \\ - 1 & \text { if } x <   0. \end{array} \right.
\]

The endomorphism \(u = s(b)\) is unitary.

\(d)\) For every \(\varepsilon > 0\) , we have \((|b| + \varepsilon)a = a(|b| + \varepsilon)\) .

\begin{enumerate}
\def\labelenumi{\alph{enumi})}
\setcounter{enumi}{4}
\item
  One has \(uau^{-1} = -a\) (check that \(s(b) = \lim_{\varepsilon \to 0} b(|b| + \varepsilon)^{-1}\) ) and \(ubu^{-1} = b\) .
\item
  One has \(upu^{-1} = q\) and \(uqu^{-1} = p\) .
\item
  We assume that the Hilbertian dimension of \(E_{p,q}\) is equal to that of \(E_{1-p,1-q}\) Then there exists a unitary endomorphism u of E such that \(upu^{-1} = q\) and \(uqu^{-1} = p\) . (The spaces \(E_{p,q}\) and \(E_{1-p,1-q}\) are orthogonal; consider the Hilbert space \(F = E_{p,q} \oplus E_{1-p,1-q}\) and apply the preceding case to the orthogonal projectors of \(F^{\circ}\) deduced from p and q by passing to subspaces.)
\item
  Suppose that \(\|p-q\|<1\) There exists a unitary endomorphism u of E such that \(upu^{-1}=q\) and \(uqu^{-1}=p\) .
\item
  Suppose that \(a = p - q\) is compact. The endomorphism \(k = qp\) induces, by passing to subspaces, a Fredholm linear map \(\widetilde{k}\) of the image of \(p\) in the image of \(q\) There exists a unitary endomorphism \(u\) of E such that \(upu^{-1} = q\) and \(uqu^{-1} = p\) if and only if the index of \(\widetilde{k}\) is zero.
\end{enumerate}

\begin{enumerate}
\def\labelenumi{\arabic{enumi})}
\setcounter{enumi}{34}
\tightlist
\item
  Let E be a Hilbert space and let u be a positive, self-adjoint partial operator on E.
\end{enumerate}

\begin{enumerate}
\def\labelenumi{\alph{enumi})}
\item
  For every \(t \in \mathbf{R}_{+}^{*}\) , we have \(u \circ \mathrm{R}(u, -t) \in \mathcal{L}(\mathrm{E})\) .
\item
  The map \(f\) of \(\mathbf{R}_{+}^{*}\) in \(\mathcal{L}(\mathrm{E})\) defined by \(f(t) = -t^{-1/2}u \circ \mathrm{R}(u, -t)\) is integrable with respect to Lebesgue measure on \(\mathbf{R}_{+}^{*}\) .
\item
  One has
\end{enumerate}

\[
\sqrt {u} = \frac {1}{\pi} \int_ {\mathbf {R} _ {+} ^ {*}} f (t) d t.
\]

\begin{enumerate}
\def\labelenumi{\arabic{enumi})}
\setcounter{enumi}{35}
\item
  Give an example of a normal partial operator on a complex Hilbert space whose essential spectrum is not closed in C.
\item
  Let E and F be complex Hilbert spaces, and let u be a closed densely defined partial operator from E into F. If u is surjective, then there exists \(v \in \mathcal{L}(\mathrm{F}; \mathrm{E})\) such that \(u \circ v = 1_{F}\) .
\item
  Let u be a partial operator in a complex Hilbert space E.
\end{enumerate}

\begin{enumerate}
\def\labelenumi{\alph{enumi})}
\item
  Assume that u is a positive symmetric operator. \(F(u)\) the Friedrichs extension of u and let q be the positive partial form associated with \(F(u)\) Show that \(\mathrm{F}(u)\) is the unique self-adjoint extension of \(u\) whose domain is contained in \(\operatorname{dom}(q)\) .
\item
  Suppose that \(u\) is a closed partial operator with dense domain. Let \(q\) the positive partial form of domain \(\operatorname{dom}(u)\) defined by \(q(x,y) = \langle u(x) \mid u(y) \rangle\) Prove that the partial self-adjoint operator representing \(q\) is \(v = u^{*}u\) .
\item
  We assume that u is symmetric and that \(u^{2}\) has dense domain. The partial operator \(u^{2}\) is symmetric and positive; the Friedrichs extension of \(u^{2}\) is equal to \(u^{*}u\) .
\end{enumerate}

\begin{enumerate}
\def\labelenumi{\arabic{enumi})}
\setcounter{enumi}{38}
\tightlist
\item
  Let \(n \geqslant 1\) an integer. We equip \(\mathbf{R}^n\) of Lebesgue measure.
\end{enumerate}

Let v be a positive function that is locally square-integrable on \(R^{n}\) . We denote \(\Delta\) the Laplacian on \(\mathcal{D}(\mathbf{R}^{n})\) or on \(\mathcal{D}'(\mathbf{R}^{n})\) , defined by

\[
\Delta (f) = - \sum_ {i = 1} ^ {n} \partial_ {i} ^ {2} f.
\]

Let D be the partial operator on \(L^{2}(\mathbf{R}^{n})\) with domain \(\mathcal{D}(\mathbf{R}^{n})\) such that

\[
\mathrm{D} (\varphi) = \Delta (\varphi + v \varphi).
\]

\begin{enumerate}
\def\labelenumi{\alph{enumi})}
\item
  The partial operator D is symmetric and positive.
\item
  Let \(u \in \text{dom}(D^{*})\) such that \(D^{*}u + u = 0\) . We have \(\Delta(u) + u = 0\) where the Laplacian of \(u \in \text{L}^{2}(\mathbf{R}^{n})\) is computed in the sense of distributions.
\item
  The distributions \(\Delta(u)\) and \((v+1)u\) are locally integrable functions, and the distribution \(-\Delta(u)\) is positive. (Use Exercise 32 of IV, p. 344.)
\item
  We have u = 0. \(|u|\) by infinitely differentiable functions \(w_{\varepsilon}\) by convolution, so that \(\Delta(w_{\varepsilon}) \geqslant 0\) and apply the previous question.) e) The partial operator D is essentially self-adjoint.
\end{enumerate}

\begin{enumerate}
\def\labelenumi{\arabic{enumi})}
\setcounter{enumi}{39}
\tightlist
\item
  Let E be a complex Hilbert space.
\end{enumerate}

\begin{enumerate}
\def\labelenumi{\alph{enumi})}
\item
  Let u be a partial normal operator on E and let \((j, |u|)\) be its polar decomposition. We have \(|u^{*}| = |u|\) and the polar decomposition of \(u^{*}\) is \((j^{*}, |u|)\) . We have \(j|u| = |u|j\) .
\item
  Let u be a partial self-adjoint operator on E and let \((j, |u|)\) be its polar decomposition. We have \(j = \mathrm{s}(u)\) , where the function s is defined by
\end{enumerate}

\[
\mathrm{s} (x) = \left\{ \begin{array}{l l} 1 & \text { if } x > 0 \\ 0 & \text { if } x = 0 \\ - 1 & \text { if } x <   0. \end{array} \right.
\]

\begin{enumerate}
\def\labelenumi{\alph{enumi})}
\setcounter{enumi}{2}
\tightlist
\item
  Let u be a partial self-adjoint operator on E and let \((j, |u|)\) be its polar decomposition. We denote by D the domain of \(|u|^{1/2}\) , and we define a partial Hermitian form q with domain D by setting
\end{enumerate}

\[
q (x, y) = \langle | u | ^ {1 / 2} (x) \mid j | u | ^ {1 / 2} (y) \rangle .
\]

The partial operator representing \(q\) coincides with \(u\) ; if \(u\) is positive, then \(q\) is the positive partial form associated with \(u\) .

\begin{enumerate}
\def\labelenumi{\alph{enumi})}
\setcounter{enumi}{3}
\tightlist
\item
  Let u and v be partial self-adjoint operators on E. Denote by \(q_{u}\) and \(q_{v}\) the partial Hermitian forms defined as in question c). Suppose that \(\text{dom}(q_{u}) \cap \text{dom}(q_{v})\) is dense in E. The sum of u and v in the sense of Hermitian forms is the partial operator \(u + v\) on E representing the partial form
\end{enumerate}

\[
q (x, y) = q _ {u} (x, y) + q _ {v} (x, y)
\]

with domain \(\mathrm{dom}(q_{u})\cap\mathrm{dom}(q_{v})\) . Show that the partial operator \(u+v\) is an extension of \(u+v\) ; if u and v are positive, then \(u+v\) is a positive partial self-adjoint operator.

\begin{enumerate}
\def\labelenumi{\alph{enumi})}
\setcounter{enumi}{4}
\item
  Give an example where the domain of \(u + v\) is not dense, but \(u + v\) is defined.
\item
  Let \(n \geqslant 1\) be an integer and \(S \subset R^{n}\) a closed set of Lebesgue measure zero. Let v be a positive measurable function on \(R^{n}\) with respect to Lebesgue measure such that v is locally integrable on \(R^{n}-S\) . The partial self-adjoint operator \(\Delta + m_{v}\) is defined.
\end{enumerate}

\begin{enumerate}
\def\labelenumi{\arabic{enumi})}
\setcounter{enumi}{40}
\tightlist
\item
  Let \(n \geqslant 1\) an integer. We denote by \(\Delta\) the Laplacian on \(\mathbf{R}^n\) , and we denote by \(u\) the reduction of \(\Delta\) to the subspace \(\mathcal{D}(\mathbf{R}^n - \{0\})\) .
\end{enumerate}

\begin{enumerate}
\def\labelenumi{\alph{enumi})}
\item
  The domain of u is dense in the Hilbert space \(E_{q}\) associated with the partial form q associated with \(\Delta\) if and only \(n \geqslant 2\) . (Identify the dual of \(E_{q}\) with a space of tempered distributions; if \(g \in E_{q}'\) vanishes on dom(u), show that g is a distribution with support contained in \(\{0\}\) ; conclude by using exercise 24 of IV, p. 338.)
\item
  The domain of \(u\) is dense in the Hilbert space \(\mathrm{E}_{\Delta}\) if and only if \(n \geqslant 4\) .
\end{enumerate}

\begin{enumerate}
\def\labelenumi{\arabic{enumi})}
\setcounter{enumi}{41}
\tightlist
\item
  Let E be a complex Hilbert space and let u be a closed partial operator with dense domain on E. We say that a partial operator v on E is compact relative to u if the domain of v contains the domain of u and if the restriction of v to \(\text{dom}(u)\) is a compact linear map from \(E_{u}\) in E.
\end{enumerate}

\begin{enumerate}
\def\labelenumi{\alph{enumi})}
\item
  Every partial operator \(v\) that is compact relative to \(u\) is bounded relative to \(u\) and its relative norm \(\| v \|_u\) is zero. (If the relative norm is not zero, construct a sequence \((x_n)\) in \(\text{dom}(u)\) converging to 0, such that \(\| v(x_n) \| = 1\) for all \(n\) and such that \((u(x_n))\) converges weakly to an element \(y\) of E; prove that \(y = 0\) and deduce that \(v\) is not relatively compact.)
\item
  Let \(\lambda \in \mathbf{C} - \mathrm{Sp}(u)\) Let \(v\) a partial operator on E whose domain contains \(\mathrm{dom}(u)\) . Then \(v\) is compact relative to \(u\) if and only if \(v \circ \mathrm{R}(u, \lambda)\) is a compact endomorphism of E.
\end{enumerate}

\begin{enumerate}
\def\labelenumi{\arabic{enumi})}
\setcounter{enumi}{42}
\tightlist
\item
  Let E be a complex Hilbert space.
\end{enumerate}

\begin{enumerate}
\def\labelenumi{\alph{enumi})}
\tightlist
\item
  Let u be a self-adjoint partial operator on E. For every \(\lambda \in R\) , we have \(\lambda \in \mathrm{Sp}_{\mathrm{e}}(u)\) if and only if there exists a sequence \((x_{n})\) in E that converges weakly to 0 in E and satisfies the conditions
\end{enumerate}

\[
\liminf _ {n \to + \infty} \| x _ {n} \| > 0, \quad \lim _ {n \to + \infty} (u (x _ {n}) - \lambda x _ {n}) = 0.
\]

\begin{enumerate}
\def\labelenumi{\alph{enumi})}
\setcounter{enumi}{1}
\item
  Let \(u\) and \(v\) of self-adjoint partial operators on E. If there exists \(\lambda \in \mathbf{C} - (\mathrm{Sp}(u) \cup \mathrm{Sp}(v))\) such that \(\mathrm{R}(u, \lambda) - \mathrm{R}(v, \lambda)\) is compact, then \(\mathrm{Sp}_{\mathrm{e}}(u) = \mathrm{Sp}_{\mathrm{e}}(v)\) .
\item
  Let u be a closed symmetric partial operator on E. Let \(u_{1}\) and \(u_{2}\) be self-adjoint extensions of u. If the space \(\operatorname{Ker}(u^{*}-i)\) is finite-dimensional, then \(\operatorname{Sp}_{\mathrm{e}}(u_{1})=\operatorname{Sp}_{\mathrm{e}}(u_{2})\) .
\item
  Let u be a self-adjoint partial operator on E and v a partial operator compact relative to u (see the preceding exercise). The partial operator \(u + v\) is self-adjoint and \(\mathrm{Sp}_{\mathrm{e}}(u + v) = \mathrm{Sp}_{\mathrm{e}}(u)\) .
\end{enumerate}

\chapter{Unitary representations}
\setcounter{exercisepar}{0}

For every vector space E, we denote by \(1_{E}\) the identity map of E. If E, F and G are vector spaces and if u: \(F \rightarrow E\) and v: \(G \rightarrow F\) are linear maps, one sometimes denotes by uv the linear map \(u \circ v\) from G to E.

For a Hilbert space E, one denotes by \(\langle x | y \rangle_{E}\) , or simply \(\langle x | y \rangle\) , the scalar product on E.

If X is a locally compact topological space and \(\mu\) a measure on X, we denote by \(\mathcal{L}^{p}(\mathrm{X},\mu)=\mathcal{L}_{\mathbf{C}}^{p}(\mathrm{X},\mu)\) and \(\mathrm{L}^{p}(\mathrm{X},\mu)=\mathrm{L}_{\mathbf{C}}^{p}(\mathrm{X},\mu)\) for \(p\in[1,+\infty]\) .

One denotes by e the identity element of a group G.

Let A be an associative algebra, E a topological vector space and \(\pi\colon\mathrm{A}\to\mathcal{L}(\mathrm{E})\) a representation of A in E. One says that \(\pi\) is non-degenerate if the set of vectors \(\pi(a)x\) for \(a\in\mathrm{A}\) and \(x\in\mathrm{E}\) is total in E.

\section{UNITARY REPRESENTATIONS}

In this section, the vector spaces are over K = R or C.

Recall that a linear representation of a group G in a K-vector space E is a homomorphism \(\varrho\) from G into the group \(\mathbf{GL}(\mathrm{E})\) of automorphisms of E (A, VIII, p. 387, def. 1). One says that E is the space of the representation \(\varrho\) , and that the dimension of E is the dimension of \(\varrho\) , also denoted \(\dim(\varrho)\) .

One can identify a representation of G in a K-vector space with a K{[}G{]}-module (loc. cit.), and one will use the corresponding terminology, for example concerning direct sums of representations, morphisms of representations or the representation on C obtained from a representation in an R-vector space by extension of scalars.

A representation \(\varrho\colon\mathbf{G}\to\mathbf{GL}(\mathbf{E})\) is said to be faithful if it is an injective homomorphism.

\pandocbounded{\includegraphics[keepaspectratio]{images/part003_3d1522a8.jpg}}

Note that this does not mean that the K{[}G{]}-module associated with \(\varrho\) is a faithful K-module (A, II, p. 28).

The character of a representation of G in a finite-dimensional K-vector space is the map \(\chi_{\varrho}\) from G to K such that \(\chi_{\varrho}(g) = \operatorname{Tr}(\varrho(g))\) for all \(g \in G\) , cf. A, VIII, p. 388.

\subsection{Reminders concerning continuous linear representations}

Let G be a topological group. One recalls that a continuous linear representation of G in a topological K-vector space E is a linear representation \(\varrho\) of G in E such that the map from G × E to E defined by \((g, x) \mapsto \varrho(g)x\) is continuous (INT, VIII, p. 128, § 2, n° 1, def. 1). This map defines an action of G on E and the image of \(\varrho\) is contained in \(\mathcal{L}(\mathrm{E})\) . One says that a continuous representation \(\varrho\) is bounded if its image is bounded in the space \(\mathcal{L}(\mathrm{E})\) endowed with the topology of bounded convergence.

Remark. --- If K = R and if \(\varrho\) is a continuous linear representation of G in an R-vector space E, then the map \(\varrho_{(\mathbf{C})}: g \mapsto 1 \otimes \varrho(g)\) is a continuous linear representation of G in \(\mathrm{E}_{(\mathbf{C})}\) ; if \(\varrho\) is bounded, then \(\varrho_{(\mathbf{C})}\) is bounded.

For every topological K-vector space E, the homomorphism which to every \(g \in G\) associates \(1_{E}\) is a continuous linear representation of G in E, called the trivial representation of G in E.

Let H be a subgroup of G. The restriction of \(\varrho\) to H is a continuous linear representation of H in E, denoted \(\mathrm{Res}_{\mathrm{H}}^{\mathrm{G}}(\varrho)\) .

Let \(\varrho_{1}\) and \(\varrho_{2}\) of continuous linear representations of G in K-vector topological spaces E \(_{1}\) and E \(_{2}\) . A G-morphism \(u\) of \(\varrho_{1}\) in \(\varrho_{2}\) is a morphism of linear representations that is continuous, that is, a continuous linear map \(u\colon\) E \(_{1} \to\) E \(_{2}\) such that \(u \circ \varrho_{1}(g) = \varrho_{2}(g) \circ u\) for all \(g \in\) G. We denote by Hom \(_{G}\) ( \(\varrho_{1}, \varrho_{2}\) ) the vector space of G-morphisms from \(\varrho_{1}\) in \(\varrho_{2}\) . It is a closed subspace of the space \(\mathcal{L}\) (E \(_{1}\) ; E \(_{2}\) ) equipped with the topology of pointwise convergence.

Let \(\varrho\) a continuous linear representation of a group G in a K-vector topological space E. The identity map \(1_{\mathrm{E}}\) of E is a G-morphism of \(\varrho\) in \(\varrho\) , which is denoted \(1_{\varrho}\) . If \(\varrho_1\) , \(\varrho_2\) and \(\varrho_3\) are continuous linear representations of G into K-topological vector spaces \(\mathrm{E}_1\) , \(\mathrm{E}_2\) , \(\mathrm{E}_3\) respectively, and if \(u: \mathrm{E}_1 \to \mathrm{E}_2\) and \(v: \mathrm{E}_2 \to \mathrm{E}_3\) are G-morphisms, then \(v \circ u\) is a G-morphism.

A G-morphism \(u\) of \(\varrho_{1}\) in \(\varrho_{2}\) is a G-isomorphism if there exists a G-morphism \(v\) of \(\varrho_{2}\) in \(\varrho_{1}\) such that \(v\circ u=1_{\varrho_{1}}\) and \(u\circ v=1_{\varrho_{2}}\) It is necessary and sufficient for this that \(u\) be a G-morphism and an isomorphism of topological vector spaces of the space \(\varrho_{1}\) in the space of \(\varrho_{2}\) ; its inverse \(v=u^{-1}\) is indeed then a G-morphism. If there exists a G-isomorphism from \(\varrho_{1}\) in \(\varrho_{2}\) we say that these representations are isomorphic.

DEFINITION 1. --- Let \(\varrho\) a continuous linear representation of a topological group G in a topological vector space E. \(\pi\) a subrepresentation of \(\varrho\) If F is a closed subspace of E and if, for every \(g \in G\) the space F is stable under \(\varrho(g)\) and \(\pi(g)\) is the endomorphism of F induced by \(\varrho(g)\) by passage to the subspaces.

Let \(\varrho\) a continuous linear representation of a topological group G in a topological vector space E. The closure of \(\{0\}\) in E is a subrepresentation of \(\varrho\) More generally, the closure of a subspace stable under the endomorphisms \(\varrho(g)\) is a subrepresentation of \(\varrho\) .

A subrepresentation \(\pi\) of \(\varrho\) is uniquely determined by a closed subspace F, stable under all endomorphisms \(\varrho(g)\) This last condition will often be stated in the form “F is a G-invariant subspace of E.” One will then also say that F defines a subrepresentation of \(\varrho\) or sometimes, by abuse of language, that F is a subrepresentation of \(\varrho\) , or of E.

Let F be a subspace of E defining a subrepresentation \(\pi\) of \(\varrho\) . For every \(g \in G\) , the linear map \(\varrho(g)\) defines, by passing to the quotient, an endomorphism \(\widetilde{\varrho}(g)\) of E/F. The map \(g \mapsto \widetilde{\varrho}(g)\) is a continuous linear representation of G in E/F, and the canonical projection of E onto E/F is a G-morphism. We say that \(\widetilde{\varrho}\) is the quotient representation of G on E/F; it is also denoted \(\varrho/\pi\) .

The subspace of elements of E invariant under the action of G is a trivial subrepresentation of \(\varrho\) . We denote it by \(\varrho^{G}\) or \(E^{G}\) .

Let A be a subset of E. The closed subspace F generated by the elements \(\varrho(g)x\) , where \(g \in G\) and \(x \in A\) is a subrepresentation of \(\varrho\) called the subrepresentation of \(\varrho\) generated by A; if F = E, one says that A generates \(\varrho\) . If the subrepresentation F is finite-dimensional, one says that the subset A is G-finite.

Suppose that A is reduced to a single element \(x \in E\) . If A generates \(\varrho\) , one says that x is a cyclic vector of \(\varrho\) and that \(\varrho\) is a cyclic representation; one says that x is a G-finite vector if A is G-finite.

The set of G-finite vectors of \(\varrho\) is a vector subspace of E stable under \(\varrho\) ; it is not necessarily closed.

Let \((\varrho_{i})_{i\in\mathrm{I}}\) a family of continuous linear representations of G in locally convex topological K-vector spaces \((\mathrm{E}_{i})_{i\in\mathrm{I}}\) . Let E be the direct sum space of the spaces \(E_{i}\) . The map \(g\mapsto(\varrho_{i}(g))\) is a linear representation of G in E, called the sum, or direct sum, of the representations \(\varrho_{i}\) If I is finite, it is continuous.

If the representations \(\varrho_{i}\) are equal to a representation \(\varrho\) for all \(i \in I\) , we say that the direct sum of the representations \(\varrho_{i}\) is the sum of Card(I) copies of the representation \(\varrho\) , and it is also denoted \(\varrho^{\mathrm{Card(I)}}\) , or Card(I) \(\varrho\) .

Let \((\varrho_{i})_{i\in\mathrm{I}}\) and \((\pi_{j})_{j\in\mathrm{J}}\) of finite families of continuous linear representations of G in topological K-vector spaces \((\mathrm{E}_{i})_{i\in\mathrm{I}}\) and \((\mathrm{F}_{j})_{j\in\mathrm{J}}\) , respectively. The canonical isomorphism of K-vector spaces

\[
\operatorname{Hom} _ {\mathrm{K}} \left(\bigoplus_ {i \in \mathrm{I}} \mathrm{E} _ {i}, \bigoplus_ {j \in \mathrm{J}} \mathrm{F} _ {j}\right)\rightarrow \bigoplus_ {(i, j) \in \mathrm{I} \times \mathrm{J}} \operatorname{Hom} _ {\mathrm{K}} (\mathrm{E} _ {i}, \mathrm{F} _ {j})
\]

(A, II, p. 13, cor. 1) induces, by passage to subspaces, an isomorphism

\[
\mathrm{Hom} _ {\mathrm{G}} \Bigl (\bigoplus_ {i \in \mathrm{I}} \varrho_ {i}, \bigoplus_ {j \in \mathrm{J}} \pi_ {j} \Bigr) \to \bigoplus_ {(i, j) \in \mathrm{I} \times \mathrm{J}} \mathrm{Hom} _ {\mathrm{G}} (\varrho_ {i}, \pi_ {j})\tag{1}
\]

which is likewise said to be canonical.

Let \(\varrho_{1}\) and \(\varrho_{2}\) of linear representations of G in K-vector spaces E \(_{1}\) and E \(_{2}\) . For \(u \in \mathcal{L}(\mathrm{E}_{1};\mathrm{E}_{2})\) and \(g \in \mathrm{G}\) we set \(\varrho(g)u = \varrho_{2}(g) \circ u \circ \varrho_{1}(g^{-1})\) . The map \(g \mapsto \varrho(g)\) is a linear representation of G in \(\mathcal{L}(\mathrm{E}_{1};\mathrm{E}_{2})\) . The space of invariant elements of this representation coincides with \(\mathrm{Hom}_{\mathrm{G}}(\varrho_{1},\varrho_{2})\) .

Let \(\varrho\) a continuous linear representation of G in a locally convex space E. Recall (INT, VIII, p. 131, § 2, no. 2) that the contragredient representation \(\breve{\varrho}\) of \(\varrho\) is the linear representation of G in the dual E' of E defined by \(\breve{\varrho}(g) = {}^{t} \varrho(g^{-1})\) .

\subsection{A criterion for continuity}

The following result makes it possible to verify that certain linear representations of a direct product of topological groups are continuous.

Lemma 1. --- Let G and H be topological groups and E a Banach space. Let \(\varrho\) be a bounded representation of G in E and \(\pi\) a bounded representation of H in E. Suppose that \(\varrho(g)\) commutes with \(\pi(h)\) for all \((g,h)\in\mathrm{G}\times\mathrm{H}\) . The map \(\varpi\) of G \(\times\) H in \(\mathbf{GL}(\mathrm{E})\) defined by \((g,h)\mapsto\varrho(g)\circ\pi(h)\) is a continuous and bounded linear representation of G \(\times\) H in E.

The map \(\varpi\) is a linear representation of G × H in E; let us verify that it is continuous. Since \(\| \varpi(g, h) \| \leqslant \| \varrho(g) \| \| \pi(h) \|\) for all \((g, h) \in \mathrm{G} \times \mathrm{H}\) , the representation \(\varpi\) is bounded, hence its image is equicontinuous in \(\mathcal{L}(\mathrm{E})\) . It then suffices to show that, for every \(x \in \mathrm{E}\) , the map \((g, h) \mapsto \varpi(g, h)x\) is continuous (remark 2 of INT, VIII, p. 129, § 2, no. 1). Let \((g_0, h_0) \in \mathrm{G} \times \mathrm{H}\) . Set \(y = \pi(h_0)x\) . For every \((g, h) \in \mathrm{G} \times \mathrm{H}\) , it follows that

\[
\begin{array}{r l} & {\| \varpi (g, h) x - \varpi (g _ {0}, h _ {0}) x \| \leqslant \| \varrho (g) (\pi (h) x - y) \| + \| \varrho (g) y - \varrho (g _ {0}) y \|} \\ & {\qquad \leqslant \| \varrho (g) \| \| \pi (h) x - y \| + \| \varrho (g) y - \varrho (g _ {0}) y \|.} \end{array}
\]

Since \(\varrho\) is bounded and that \(\varrho\) and \(\pi\) are continuous, this implies the assertion.